\documentclass[10pt,a4paper,twocolumn]{article}
\usepackage[top=21mm,bottom=22mm,left=19mm,right=19mm,columnsep=8mm]{geometry}
\usepackage{array,longtable}
\usepackage[T1]{fontenc}
\usepackage[english]{babel}
\usepackage{mathptmx}
% Removed missing external fontsetup.tex; Times-compatible mathptmx is retained.
\usepackage[final]{microtype}
\usepackage{balance}
\usepackage{amsmath,amssymb,amsthm}
\allowdisplaybreaks[1]
\usepackage{booktabs,graphicx,url}
\usepackage[hidelinks]{hyperref}
\setlength{\emergencystretch}{1.3em}
\newtheorem{theorem}{Theorem}
\newtheorem{proposition}{Proposition}
\newtheorem{remark}{Remark}
\newtheorem{lemma}{Lemma}
\newcommand{\C}{\mathcal C}
\newcommand{\U}{\mathcal U}
\newcommand{\HH}{\mathcal H}
\newcommand{\norm}[1]{\left\|#1\right\|}

\hypersetup{pdftitle={A Decentralized Monetary Economy with Delayed Production, Employment, and Education},pdfauthor={Jan Olsina}}

\begin{document}
\title{A Decentralized Monetary Economy with Delayed Production, Employment, and Education}
\author{Jan~Ol\v{s}ina\thanks{Independent Researcher, Prague, Czech Republic.
Email: \href{mailto:jan.olsina@gmail.com}{jan.olsina@gmail.com}.
ORCID: \href{https://orcid.org/0009-0003-9816-3958}{0009-0003-9816-3958}.}}
\date{26 October 2026}
\maketitle

\begin{abstract}
We specify a discrete-time, spatial agent-based monetary economy intended
as a reproducible foundation for an economic simulation game. Households
value nutritional, residential, material, and social services with
nonhomothetic satiating preferences. Firms transform dated inputs and
skill-specific labor into outputs through capacity-constrained recipes.
Local markets adjust positive price quotes using funded orders and
committed supply; agents possess neither the global economic state nor a
counterfactual economy simulator. Employment, housing, and education are
allocated through local contracts and matching, with bounded preferences
for occupational and environmental attributes. Explicit settlement,
reservation, and production rules prevent double spending and premature
output. We establish accounting feasibility, conservation of money in the
closed reference economy, information locality, and nominal scale
equivariance. A local linear stability criterion and a separate convex
benchmark identify the limited settings in which convergence can be
analyzed. No global convergence or empirical calibration of the complete
economy is asserted. Appendices provide a 34-good industrial catalog,
production coefficients, utility parameters, and building characteristics.
\end{abstract}

\section{Introduction and scope}
The desired simulation combines heterogeneous household needs, material
production chains, spatial access, employment, and education without
assigning intrinsic monetary values to commodities. Prices must be
generated by transactions and local adjustment, while preferences and
technologies remain physical or behavioral primitives. The objective is a
precise disequilibrium mechanism, not an algorithm that presupposes its
own equilibrium solution.

Decentralized exchange and the emergence of prices have been studied in
agent-based general-equilibrium models \cite{gintis2007}. The
Eurace@Unibi program connects spatial production, consumption, labor, and
human-capital mechanisms \cite{dawid2014}. These precedents motivate the
architecture, but their numerical findings do not establish properties of
the present rules. Classical instability examples also show why an
excess-demand price rule cannot be assumed to converge for arbitrary
preferences \cite{scarf1960}. Recent work studies price adjustment with
forward-looking behavior and explicit conditions for stability
\cite{werning2026}; those assumptions differ from the adaptive information
structure used here.

This article defines a reference algorithm and distinguishes exact
invariants from conditional analytical results. The industry appendix is
a hypothetical game economy, not an engineering calibration. Public
finance is included as a cash-constrained institutional actor. Banking,
unsecured credit, endogenous crime, births, deaths, migration, and an
unmodeled external buyer are excluded. Initial assets, knowledge, monetary
balances, and physical reserves are exogenous endowments. They are never
automatically replenished.

\subsection{Purpose, design concepts and documentation}
This is a model-specification article organized with an ODD crosswalk
\cite{grimm2020}. Purpose and intended patterns are in this introduction;
entities, state and scheduling are in Sec.~2; adaptation, objectives,
sensing, interaction and learning are specified in Secs.~3--6.
Emergent outcomes are prices, employment, consumption and skill allocation.
Stochasticity consists of fixed taste draws, review offsets, tie priorities,
illness and injury. Households are collectives with explicit bargaining
weights. Initialization, input data and observations are in Sec.~9 and
the appendices. No fitted empirical pattern or full dynamic validation
is claimed; the static check tests only resource feasibility.

\section{State, chronology, and information}
\subsection{Entities and units}
Let $t\in\{0,1,\ldots\}$ index days. Let $\HH$, $\mathcal I$,
$\mathcal F$, $\mathcal L$, and $\mathcal G$ denote households,
persons, establishments, locations, and commodity classes. Each person
belongs to one household. Establishments include factories, farms,
landlords, schools, retailers, service providers, and transport operators.
The government is an additional account with an explicit policy rule.
Age in days increases by one each day. Children aged 6--17 may attend
basic education; adult remedial students are also eligible. Persons below
18 cannot accept paid employment. At 18 they retain their household,
assets and attained skills and become eligible for work and adult courses.
Consumption equivalents are 0.5 below age 6, 0.75 from 6 to 17, and 1
thereafter; clothing and occupancy still count individual persons.
School occupies five of a child's fourteen available daily hours;
there is no child productive labor. Births, deaths and migration are
disabled, so the population is closed but its age composition changes.
No additional age productivity effect or retirement is modeled.

A commodity market $m=(\ell,g,q)$ identifies a location, a good, and a
verifiable grade $q$, including any relevant freshness guarantee. Goods
within one market are interchangeable for settlement. Different grades
may have different quotes. Let $p_{m,t}>0$ be the posted quote. Money is
the unit of account; no commodity price is fixed as a numeraire.
A market is a clearing institution, not a shop building. A shop owns
inventory and submits supply into the market, like any other seller;
its wholesale purchases and retail sales occur in different dated rounds.
The shop provides optional staffed access with a separately posted
handling-service fee; direct pickup from a known producer remains possible.
The handling fee has its own funded-demand/capacity update and covers
shop payroll and paper. Same-grade goods in one market retain one common
goods quote; handling and freight charges are separately itemized.

Most material quantities are divisible lots. Jobs, homes, people, school
places, and production slots are discrete. The baseline therefore includes
indivisibilities, even though a fixed-choice consumption subproblem may
be convex. Electricity and dated service capacity are flows, not
inventories. Appendix~\ref{app:goods} specifies units and storage rules.

\subsection{Complete state}
The simulator state $X_t$ contains all variables needed for a Markov
transition: cash accounts; cash and inventory reservations; inventories
by good, owner, location, and age; durable condition; assets and ownership;
jobs and leases; skills, course progress, health, and service stocks;
production and construction pipelines; transport queues; natural stocks;
price quotes; local beliefs; search sets; and contract histories.

An agent $a$ observes a filtration $\mathcal J_{a,t}$ generated by its
own state, its own transactions, received messages, visited places, known
quotes, and local market announcements. A market sees only the orders,
reservations, and settlement records needed for its own clearing. Agents
may observe the total orders and offered supply in markets they attend;
this is local information, not the state of the economy.

Policies satisfy
\begin{equation}
 A_{a,t}=\pi_a(\mathcal J_{a,t},\xi_{a,t}),
 \label{eq:policy}
\end{equation}
where $\xi_{a,t}$ is a seeded private or matching random variable. A public
research dashboard may inspect $X_t$, but its aggregates are not
automatically passed to $\pi_a$. Comparing an agent's own bundles, homes,
or job offers is permitted; evaluating alternate economy-wide equilibria
is not.

\subsection{Reference daily event order}
The following order defines the baseline transition, including its
otherwise economically consequential timing conventions.
\begin{enumerate}
\item Age preexisting lots and remove expired ones to a waste ledger.
Then release eligible production from the previous day as new age-zero
lots and complete freight. Freight retains the original production age;
delivery never refreshes a lot. Apply
contractually due transfers that can be funded.
\item On staggered review dates, update local search, job offers, leases,
and school applications. Accepted changes take effect on the next day.
Every person has one feasible daily activity schedule.
\item For today's existing contracts, reserve payroll, rent, and tuition.
Unfunded work or paid service is suspended before it is performed.
Reserve the physical inputs and staff needed to offer current utility
and service capacity.
\item Households and firms submit jointly affordable orders. Reserve
their cash and sellers' physical offers across all markets. Clear each
local market once, using the frozen beginning-of-round information.
\item Execute feasible work, transport, service appointments, and
production progress. Start new batches only with secured material
inputs and current first-stage labor capacity. Release earned payroll
and payment for performed services from escrow.
\item Households consume available goods, use durables, and update
health and social states. Update skills, condition, forecasts, and
financial accounts. Release unused reservations; update tomorrow's
quotes from the submitted orders recorded before rationing.
\end{enumerate}
Newly received market revenue cannot finance another order in the same
clearing round. End-of-day wages become spendable on the next day. This
is a working-capital convention, not a universal fact about money.
An implementation with more rounds must retain reservations and use an
adjustment rate consistent with its round length.

Daily quantities enforce time and aggregate capacity, not a
minute-by-minute transport simulation. Service sessions have explicit
time windows; incompatible appointments cannot both be accepted.
Employed workers search on a staggered seven-day cycle; unemployed or
newly suspended workers search daily. Households review their residence
and education on staggered 30-day cycles, with daily housing search when
unhoused. Firm offer review occurs with its seven-day production review.
The reference production lag is at least one daily boundary for material
outputs. Kitchen output becomes
available at the next boundary, fresh for that following day. Electricity
and other contemporaneous services use reserved delivery capacity instead.

\section{Household services and preferences}
\subsection{Satiation and material attributes}
For $z\geq0$, define
\begin{equation}
 H(z)=\frac{3\min\{z,1/2\}
 +\min\{(z-1/2)_+,1/2\}}{2}.
 \label{eq:H}
\end{equation}
It is increasing, concave, and constant after $z=1$. The first half of a
need has three times the marginal utility of the second half.
For a consumption vector $x_i$, nutritional and sensory services are
\begin{equation}
 (N_i,P_i,M_i,Q_i)
 =\sum_g (n_g,\pi_g,\mu_g,\zeta_{ig}q_g)x_{ig},
 \label{eq:nutrients}
\end{equation}
with $\zeta_{ig}\in[0.85,1.15]$ a stable taste multiplier. Nutrition is
unaffected by taste. Daily targets are $N_i=1$, $P_i=0.5$, and $M_i=0.5$
per adult; planned consumption obeys $N_i\leq1.10$. Appendix tables give
the attributes and preparation requirements.

The person utility is
\begin{align}
 u_{i,t}={}&\sum_k w_{ik}H(z_{ik,t})
 +10\log(1+L_{i,t}/2)\nonumber\\
 &+30\log\frac{h_{i,t}+5}{105}
 -C^{\rm work}_{i,t}
 +P^{\rm home}_{i,t}+P^{\rm visit}_{i,t}.
 \label{eq:utility}
\end{align}
Here $h\in[0,100]$ is health and $L\geq0$ is discretionary leisure.
The household objective uses explicit nonnegative bargaining weights
$\omega_{ih}$, with $\sum_{i\in h}\omega_{ih}=1$. Equal weights are the
baseline. Child consumption-equivalence factors alter physical targets,
not interpersonal moral worth.

Each person selects at most one clothing set per day. A set with positive
condition gives one unit of coverage; condition measures remaining life,
not the fraction of the body covered. Fine clothing additionally supplies
style proportional to condition. Only the selected set wears; at zero
condition it is retired. Furnishing and housing service can remain
condition-weighted.

Home hygiene is the minimum of normalized washing water, soap, and
facility access. A purchased bath adds delivered hygiene service; its
provider has already purchased the water and soap. The customer must not
be charged the same physical inputs a second time. Medical care affects
expected health transitions rather than adding utility in proportion to
money spent.

For $k$ in social connection, recreation, reading, and religious practice,
\begin{equation}
 s_{ik,t+1}=\frac67s_{ik,t}+\sum_v q_{v,t}a_{kv}y_{iv,t}.
 \label{eq:stocks}
\end{equation}
Here $y_{iv,t}$ is visits for a paid venue (one hospital unit is a
patient-day), and hours for an unpaid activity. Coefficients $a_{kv}$ are
service per corresponding unit; visit duration separately enters the time
budget. Library sessions are one hour, so their input coefficients per
user-hour coincide with per-session coefficients.
Stocks above one are retained and decay normally, although their current
utility is satiated. A repeated book has declining novelty. Simultaneous
borrowing of the same physical copy is prohibited.

The working-age time constraint is
\begin{equation}
 L_i+h_i^{\rm paid}+h_i^{\rm study}+h_i^{\rm travel}
 +h_i^{\rm domestic}+h_i^{\rm visits}\leq14.
 \label{eq:time}
\end{equation}
Eight additional hours are reserved for sleep and two for eating and
unavoidable personal care. Queues and searching use actual time. The
model cannot satisfy thirty daily food requirements by consuming a
monthly quantity on the first day.

\subsection{Small, stable attribute preferences}
For an option $o$ with $d$ normalized attributes $a_o\in[-1,1]^d$, set
\begin{equation}
 P_i(o)=2\tanh\left(\frac{\theta_i^\top a_o}{\sqrt d}\right),
 \qquad \theta_{ik}\sim {\rm Uniform}[-1,1].
 \label{eq:tastes}
\end{equation}
Draw separate vectors for occupations, homes, and venues once at
initialization. Features include outdoor activity, craftsmanship,
helping others, quietness, garden access, density, and familiarity.
The total venue-match term is clipped to $[-2,2]$ per day after weighting
each visit by $\min\{1,\text{visit hours}\}$. Reentering a preferred venue
therefore cannot generate unlimited utility.

Rent, wages, crowding, pollution, commuting, and physical safety are
handled in budgets, services, time, and work costs. They must not be
hidden a second time in the small preference bonus. Interest and skill
are separate variables. Multiplying a decision maker's entire objective
by a common positive number changes no choice. Within a household,
rescaling only one member's utility requires an inverse adjustment of
that member's bargaining weight to preserve collective preferences.
Changing only a taste coefficient does change preferences.

\subsection{Local household decision problem}
Let $\C_{h,t}$ contain the currently known feasible combinations of
home, job, course, and appointment choices. Search is bounded: retain up
to five known nearby suppliers per category and discover up to two new
local listings weekly, sampled from a public local directory. Never
infer an unobserved vacancy or supplier from the global state.

For each $c\in\C_{h,t}$, household $h$ solves its own finite-horizon
problem
\begin{equation}
 \max_{a_0,\ldots,a_{T-1}}
 \sum_{\tau=0}^{T-1}\beta^\tau
 \widehat{\U}_{h,t+\tau}(a_\tau\mid c)
 +\beta^T\Psi_{h,t}(X_{h,T}),
 \label{eq:hhplan}
\end{equation}
subject to inventory, time, qualification, occupancy, and cash
constraints. $\U_h=\sum_{i\in h}\omega_{ih}u_i$ and $\beta=0.9999$.
Use $T=30$ for routine shopping, $T=90$ when evaluating a move or job
change, and $T$ equal to remaining course duration plus 360 days when
evaluating education. These are local planning horizons, not perfect
foresight. Execute only the current-day decision and replan tomorrow.

Forecast quotes are the latest observed quotes, held constant. Existing
contract incomes are forecast at observed payroll reliability; prospective
post-education income uses observed qualified vacancies and the local
employment frequency described below. Purchases are forecast at the
agent's observed fill fraction. All actual executions use realized fills,
not these forecasts. Expected future earnings are never current cash.
The baseline evaluates utility along these forecast mean paths: it is a
certainty-equivalent planning rule, not an assertion that utility of a
mean outcome equals expected utility. Any risk-sensitive alternative
must replace this forecast convention explicitly.

For specificity, let the terminal continuation be
\begin{equation}
\begin{aligned}
 \Psi_{h,t}={}&\lambda_{h,t}W^*_{h,t}
 \log(1+W_{h,T}/W^*_{h,t}),\\
 W_{h,T}={}&m_{h,T}+\sum_g\chi_g\hat p_{hg,t}I_{hg,T}.
\end{aligned}
\label{eq:terminal}
\end{equation}
where $W^*_{h,t}=30c^{\rm ess}_{h,t}>0$ and usable remaining durables enter condition-weighted $I$, expired
goods have value zero, and $\chi_g=1/2$ is a conservative resale belief.
Only goods with known resale outlets enter the sum. This is a private
terminal approximation, not a guaranteed redemption facility or an
intrinsic price. Let $V_{h,t}(b)$ be the household's one-day attainable
utility at current known offers with its discrete arrangements fixed.
Define
\begin{equation}
 \lambda_{h,t}=
 \frac{V_{h,t}(B_{h,t}+\delta b)-V_{h,t}(B_{h,t})}{\delta b},
 \quad \delta b=0.01\,c_{h,t}^{\rm ess}.
 \label{eq:lambda}
\end{equation}
Here $c^{\rm ess}>0$ is the cost of a fixed physical reference basket at
known positive quotes, including emergency last-observed quotes when a
seller is temporarily unavailable. This basket defines a local monetary
scale only; it does not create an order or fix prices. $\lambda$ may be
zero when all attainable one-day services are satiated. Terminal
approximations can affect saving and are part of the model's sensitivity
analysis.

Today's discretionary budget, after funded fixed obligations, is
\begin{equation}
 B_{h,t}=\min\{m^{\rm free}_{h,t},
 [\hat y_{h,t}+\kappa_h(m^{\rm free}_{h,t}-m^*_{h,t})]_+\},
 \label{eq:budget}
\end{equation}
where $\kappa_h=0.1$, $m^*_{h,t}=7\hat e^{\rm ess}_{h,t}$, and income
and essential expenditures are exponential averages with gain $0.1$.
The same rule is used in local forecast periods. Initial histories are
scenario data. Purchases must fit current free cash even when
Eq.~\eqref{eq:hhplan} forecasts a profitable investment.

Equation~\eqref{eq:hhplan} defines a preference criterion. Its operational
policy uses bounded candidate evaluation, not exhaustive dynamic search.
Daily shopping fixes home, jobs, course and durable assignments and solves
the one-day continuous problem with the concave terminal-wealth term.
For each scheduled discrete review, evaluate the status quo and at most
ten alternatives, changing only one arrangement or one member's job.
For each candidate, simulate the stated horizon by repeating this one-day
policy using frozen quotes and forecast fills, updating stocks, cash,
course progress and contract expiry. Choose the best discounted rollout.
This is a stationary-policy rollout, not an exact finite-horizon optimum.
Use piecewise-linear interpolation of logarithms on a fixed dimensionless
grid of 65 equally spaced points over the feasible argument interval;
all other satiation pieces are exact. The daily fixed-choice problem is
an LP when health effects are held at their current forecast.
Long-horizon rollouts occur only on review dates. Ties prefer the current
arrangement, then lower spending, then stable identifiers. Report LP count,
runtime and objective approximation error in computational experiments.

\begin{proposition}[Restricted consumption concavity]
Fix discrete choices, activity schedules, wear assignments, and expected
health effects. Suppose remaining service transitions are affine or
concave nonnegative minima of affine resources, and use divisible
quantities. Then the one-period continuous utility maximization problem
with linear resource constraints is a convex optimization problem
(maximization of a concave function). With a piecewise-linear leisure
approximation it admits a linear-program representation.
\end{proposition}
\begin{proof}
Equation~\eqref{eq:H} is nondecreasing and concave; composition with each
stated concave service map preserves concavity. Leisure utility is
concave in an affine nonnegative residual. Fixed match terms are
constants and linear resource constraints define a convex set. For the
linear-program form, use two service segments, each bounded by $1/2$,
with slopes $3w_k/2$ and $w_k/2$ and total segment use bounded by the
available service. Positive decreasing slopes fill the first segment
before the second. Standard convex-optimization conventions apply
\cite{boyd2004}.
\end{proof}
Discrete investments, illness probabilities, and the complete dynamic
model are not covered by this proposition.

\section{Production, procurement, and investment}
\subsection{Physical production pipelines}
A recipe $r$ has material input vector $a_r$, output vector $b_r$,
skill-specific effective labor requirements $L_{kr}$, electricity
requirement $E_r$, and minimum elapsed duration $\tau_r\geq1$.
Starting batch $j$ immediately locks and consumes $a_r$ and occupies one
compatible slot. Let $v_{j,t}\in[0,1]$ be accumulated progress. Its daily
increment is
\begin{equation}
 \Delta v_{j,t}=\min\left\{1-v_{j,t},\tau_r^{-1},
 \min_{k:L_{kr}>0}\frac{H_{jk,t}}{L_{kr}},
 \frac{e_{j,t}}{E_r}\right\},
 \label{eq:progress}
\end{equation}
omitting the last ratio when $E_r=0$. $H_{jk,t}$ and $e_{j,t}$ are
allocated effective hours and delivered electricity. Only resources
corresponding to the chosen increment are used; unused allocations may
serve another task. The per-day cap preserves minimum process duration.
Output $b_r$ becomes available at the next boundary after $v_j=1$.
Materials are never purchased or consumed a second time after a pause.

Required site resources, installed equipment, and worker qualifications
are additional feasibility conditions. Within a firm allocate daily
capacity first to contractual service obligations, then to unfinished
batches by oldest start date, then to new batches by recipe identifier.
Workers' actual hours are disjoint across tasks. General labor contributes
effective hours according to skill productivity, but paid hours and time
cost remain actual hours. Productivity is defined by Eq.~\eqref{eq:productivity}.

Services and electricity cannot be sold from future inventory. A provider
offers only capacity backed by current inputs, reserved payroll, and
feasible schedules. Consumers pay into escrow; providers receive payment
only for performed units. The reference power network has no storage and
enforces daily delivered energy capacity; finer network dispatch is an
extension. A non-electric well and initial water/fuel stocks permit
bootstrapping the water--electricity cycle.

\subsection{Contemporaneous power capacity}
Power is a flow activity, not a delayed recipe. Let $z_{f,t}$ be operating
ten-unit blocks, continuously divisible. With installed condition $c_f$,
\[
\begin{aligned}
 0&\le z_{f,t}\le 12c_f,\\
2z_{f,t}&\le I^{\rm charcoal}_f,\\
0.2z_{f,t}&\le I^{\rm water}_f,\\
(4/3)z_{f,t}&\le H^{\rm technical}_f,\\
(8/3)z_{f,t}&\le H^{\rm general}_f.
\end{aligned}
\]
Delivered generation is $10z_{f,t}$. The six-worker cap and actual-hour
schedules apply separately. Reserve inputs and payroll before offering
capacity; choose realized $z$ equal to funded delivered demand divided by
ten, debit proportional inputs, and release unused reservations.
Unserved electricity is never banked. Non-electric wells instead produce
one storable water lot per four actual hours, released at the next
boundary, with source capacity specified by the scenario.

\subsection{Sales beliefs and output release}
Each firm updates local forecasts with gain $\alpha=0.1$:
\begin{equation}
 \hat z_{f,t+1}=(1-\alpha)\hat z_{f,t}+\alpha z^{\rm obs}_{f,t}.
 \label{eq:forecast}
\end{equation}
For demand, $z^{\rm obs}$ includes unfilled funded orders, not merely
sales. In a common market, assign attempted orders for forecasting by
\begin{equation}
 D_{fm,t}=\sigma_{fm,t}D_{m,t},\qquad
 \sigma_{fm,t}=
 \frac{q^{\rm sold}_{fm,t-1}+\nu C_{fm}}
 {\sum_{j\in m}(q^{\rm sold}_{jm,t-1}+\nu C_{jm})},
 \label{eq:share}
\end{equation}
where $C_{fm}$ is the registered physical daily output capacity and
$\nu=0.05$. This is an observable local demand allocation used for
forecasting, not a priority right in actual clearing. If the denominator
is zero, no firm is assigned an order. Entry requires an actual
advertised technology and financing.

Let $I^{\rm sell}_{fg}$ be saleable output stock after existing
commitments. The firm keeps a target buffer $I^*_{fg}=b_g\hat d_{fg}$,
where $b_g=\min\{3,\text{shelf life}/2\}$ days, with $b_g=3$ for
nonperishables. Its offered quantity is
\begin{equation}
 S_{fg,t}=\min\{I^{\rm sell}_{fg,t},
 [\hat d_{fg,t}+\gamma(I^{\rm sell}_{fg,t}-I^*_{fg,t})]_+\},
 \label{eq:release}
\end{equation}
with $\gamma=0.25$. Lots expiring at the next boundary are offered in
full, provided they meet the market's stated grade. A plant need not dump
its complete buffer each day; otherwise an inventory target can create
permanent artificial excess supply.

\subsection{A specified local planning policy}
On its staggered seven-day review a firm uses
\begin{equation}
\begin{aligned}
 T_f=\max\bigl\{&31,\; 7+\max_{r\in f}\tau_r,\\
 &1+\text{longest remaining contract term}\bigr\}.
\end{aligned}
\end{equation}
An accepted offer includes one paid onboarding day and thirty productive
days, hence all thirty-one days enter the hiring comparison. Existing
workers beyond expiry are forecast available only under an explicit
renewal decision at the current offer, with their observed acceptance
frequency; an infeasible renewal generates no labor.

Plans use a time-indexed mixed-integer linear program (MILP): binary
batch-start variables for each slot and day; continuous progress, input
orders, inventories, output sales, cash and labor allocations; and binary
completion indicators. Progress obeys Eq.~\eqref{eq:progress} as linear
upper bounds with input debits at start, completion allowed only after
unit accumulated progress, and output released one boundary later.
Finite slot, storage and cash bounds provide finite indicator bounds.
Inputs are separated into age cohorts, expired lots are removed, and
sales are bounded by forecast demand and available output. Actual-hour
labor is bounded by funded contracts, with fixed known skill multipliers.
Future quotes and fill fractions are frozen private beliefs.

The objective for both operation and investment is
\begin{equation}
\begin{aligned}
 J_f={}&\sum_{\tau=0}^{T_f-1}\beta^\tau
       \bigl(\widehat{\mathrm{receipts}}_\tau\\
       &\quad-\widehat{\mathrm{payments}}_\tau\bigr)
       +\beta^{T_f}V_f^{\rm terminal},\\
 &\beta=0.9999.
\end{aligned}
\label{eq:firmplan}
\end{equation}
Cash feasibility is enforced separately each day; terminal cash is not
added again. Terminal inventories use half the last known resale quote.
For unfinished batch $j$ with progress $v_j$, assign
\[
 V_j^{\rm WIP}=\left[
 \tfrac12\sum_g\hat p_g b_{rg}
 -(1-v_j)\widehat C_r^{\rm remaining}\right]_+ .
\]
Remaining cost includes remaining labor and electricity, but excludes
already consumed inputs. Assign zero if no known feasible completion
route or resale outlet exists. This is a private accounting approximation,
not a liquidation guarantee. To avoid speculative terminal accumulation,
new starts must have a forecast sales outlet within one recipe duration
after the horizon, and total planned sales plus terminal output cannot
exceed demand over that extended interval. Existing capital salvage is
condition-weighted at known resale quotes; no outlet means zero.
Investment uses the identical objective with a 360-day horizon.

Terminal positive-part valuations use bounded indicator variables in the MILP.
The MILP search stops after 10,000 branch-and-bound nodes or a relative
gap of 1\%, whichever occurs first. Execute only a feasible incumbent;
the all-idle plan with already committed bills is the fallback if feasible,
otherwise suspend unfunded activities. Record the incumbent, bound,
relative gap, solver version and deterministic settings. No optimality
or game-scale runtime is claimed for this bounded algorithm.

Let $R_{fg}$ be the material requirement of the selected near-term plan,
$J_{fg}$ committed inbound deliveries, and $I^{\rm input}_{fg}$ usable
input stock. Procurement targets are
\begin{equation}
 q^{\rm target}_{fg}=[R_{fg}-I^{\rm input}_{fg}-J_{fg}]_+.
 \label{eq:procure}
\end{equation}
Define the procurement envelope after today's escrow reservations by
\[
 B^{\rm proc}_{f,t}=
 [m^{\rm free}_{f,t}-\widehat K^{\rm fixed}_{f,t+1:t+7}]_+,
\]
where $\widehat K^{\rm fixed}$ is the undiscounted payroll, rent, tuition,
contractual debt-free service bills and essential maintenance over the
next seven days, excluding today's already reserved amounts and all
discretionary input purchases. Expected receipts do not enlarge this
cash envelope. Orders obey $\sum_g\hat p^{\rm delivered}_{fg}q_{fg}
\leq B^{\rm proc}_{f,t}$.
When the complete target is unaffordable, allocate the envelope by
progressive filling of input coverage
\begin{equation}
 c_{fg}=\frac{I^{\rm input}_{fg}+J_{fg}+q_{fg}}{R_{fg}}
 \quad(R_{fg}>0).
\end{equation}
Raise the smallest coverage first, up to one, using known delivered
quotes; ties split expenditure to equalize coverage. This is a finite
water-filling procedure, not independent spending of the entire budget
on every input. It does not guarantee that independent market rationing
delivers a complete bundle. Incomplete bundles remain real inventories
and affect the next plan. Canceling a project does not refund consumed
inputs or undo completed trades.

Each plan implies required worker-hours by skill. Firms announce the
unfilled portion as vacancies, subject to their maximum profitable wage
and ability to reserve the next day's payroll. For the maximum wage, compare a frozen feasible without-hire incumbent
with a with-hire plan over the same horizon, including all thirty-one
paid days and onboarding. Search a wage grid with spacing
$0.01c^{\rm ess}_{\rm adult,t}/8$ per hour, up to the no-wage gross-receipt
plus salvage bound divided by discounted committed hours. Choose the
largest grid wage whose feasible with-hire incumbent has nonnegative
incremental discounted value and a feasible cash path.
This defines a conservative, solver-dependent reservation wage, not an
exact global profit threshold. Reserve next-day payroll before acceptance.
Subsequent inability to pay invokes the explicit suspension rule; the
model does not describe these contracts as fully prepaid guarantees.

Every 30 days an owner may consider one additional building module from
its known catalog. For each candidate, evaluate a 360-day cash forecast
at observed quotes, include construction time and all bills of materials,
and add the same half-quote terminal salvage convention. Start only if
incremental discounted value is positive and initial commitments can be
funded without consuming reserved payroll. Choose the largest positive
value; ties use stable module identifiers. The default discount factor
is $0.9999$ per day. A household may found a firm by the same rule using
its unreserved savings. A 30-day idle firm with insufficient maintenance
funds enters liquidation, offering inventories and reusable assets in
their actual markets. Sale receipts belong to its owners after funded
contractual obligations. No bailout is automatic.

\section{Local exchange and price formation}
\subsection{Reservations and actual settlement}
Let $m^{\rm free}_{a,t}$ be cash after preexisting reservations. A buyer's
orders across all attended markets satisfy
\begin{equation}
 \sum_m p_{m,t}d_{am,t}
 +\text{reserved delivery charges}
 \leq m^{\rm free}_{a,t}.
 \label{eq:escrow}
\end{equation}
Household discretionary orders also satisfy Eq.~\eqref{eq:budget}.
Sellers reserve particular inventory lots or dated capacities. The
reservation ledger is account-local, allowing atomic rejection of a
duplicate commitment without revealing the rest of the economy.

For a divisible material market, define
\begin{equation}
 D_m=\sum_a d_{am},\qquad S_m=\sum_f s_{fm},\qquad
 Q_m=\min\{D_m,S_m\}.
\end{equation}
Executed quantities are
\begin{align}
 x_{am}&=d_{am}\min\{1,S_m/D_m\},\label{eq:buyfill}\\
 y_{fm}&=s_{fm}\min\{1,D_m/S_m\},\label{eq:sellfill}
\end{align}
with zero execution if the relevant total is zero. All trades in the
round settle at the posted $p_m$. Transfer $p_mQ_m$ from buyers to
sellers, transfer ownership of $Q_m$ goods, and release unused cash and
inventory reservations. Select seller lots by earliest expiry, then lot
identifier. Goods at another location require separately funded freight;
they are unavailable for use until delivered.

Nondivisible homes, jobs, and appointments use the matching procedures
below; their allocation must not be implemented as fractions of persons.
Services that fail to be delivered release the corresponding escrow to
the customer. Failed transport leaves goods with their current owner at
the documented location; it does not duplicate them.

Households face the same complementary-input problem as firms. Orders
are formed from complete intended cooking bundles net of usable stocks,
but proportional fills need not preserve these bundles. After clearing,
solve the fixed-choice consumption LP again against actual stocks with
no further purchases; consume only recipes whose full food, fuel, water
and time inputs are available. Preserve unmatched inputs until used or
expired, and record unmet nutrition and stranded inventory. Ready-to-eat
alternatives compete in the ex-ante shopping problem. The mechanism
guarantees feasible consumption, not satisfaction of planned nutrition.

\subsection{Price update}
At the end of the round, record funded orders \emph{before rationing} and
committed supply \emph{before execution}. Let
\begin{align}
 e_{m,t}&=\frac{D_{m,t}-S_{m,t}}
 {D_{m,t}+S_{m,t}+\varepsilon_m},\label{eq:imbalance}\\
 p_{m,t+1}&=p_{m,t}\exp(\eta_m e_{m,t}).
 \label{eq:price}
\end{align}
Set $\varepsilon_m=0.01$ of the market's physical trading unit and
$\eta_m=0.02$ per daily round in the reference parameterization.
The numerical coefficients bound daily log-price changes, so 0.02 means
at most about 2.02\% per round, not a fixed 2\% change.
These are quantity regularization and adjustment-speed parameters, not
price anchors. If $D_m=S_m=0$, the quote is unchanged.

Actual purchases always equal actual sales; substituting either for
$D_m$ and $S_m$ would erase the imbalance signal. Previous unfilled
orders expire and must be funded again before they count as new demand.
This is a posted-price non-t\^atonnement process: trades occur away from
equilibrium and supply is committed within each round. It belongs to the
tradition studied by Hahn and Negishi \cite{hahn1962}, but their specific
stability result is not imported into this monetary rationing model.
The rule is a market institution. It is neither a demonstrated
profit-maximizing seller policy nor a strategy-proof mechanism. Strategic
withholding and order manipulation are outside the price-taking baseline.

\begin{remark}[Zero prices and nonexistent supply]
Equation~\eqref{eq:price} preserves strictly positive finite quotes over
every finite horizon. A free-disposal good can approach a zero price
asymptotically; disposal itself may require a paid service. A positive
quote does not establish that a good is produced or available. Entry,
financing, and physical feasibility remain necessary.
\end{remark}

\section{Employment, residence, and education}
\subsection{Jobs and wages}
For a job $j$ requiring $h_{ij}$ actual hours, use
\begin{equation}
 C^{\rm work}_{ij}=\frac{h_{ij}}8
 (10e_{ij}+6\rho_i r_j-2a_j-P_i^{\rm job}(j)),
 \label{eq:workcost}
\end{equation}
with effort $e$, hazard $r$, autonomy $a$, and risk aversion
$\rho_i\in[0.7,1.3]$. Job income affects the feasible household plan,
not a separate fixed utility per currency unit. Commute time enters
Eq.~\eqref{eq:time}; travel fees enter its budget. Unemployment and
feasible self-employment are alternatives.

At each hiring review, process household members in a seeded rotating
order. At most one member of any household participates in a matching
subround; all other members' jobs are held fixed while computing that
member's preference list. Run worker-proposing deferred acceptance on
these frozen lists, with each firm ranking eligible workers by effective
productivity per paid hour and stable identifier within fixed vacancy
classes. Commit that subround before reviewing another member, and
process each member at most once per review. Previously committed
contracts cannot be displaced. Switching costs 3 utility units.
An offer accepted on day $t$ starts paid onboarding on $t+1$, with zero
productive output; work starts on $t+2$, and the last paid day is $t+31$.
The old job is released at the start of $t+1$.

Deferred acceptance terminates on each finite frozen graph
\cite{galeshapley1962}. Its subround stability does not imply stability
against joint household deviations: couples create linked preferences
\cite{roth1984}. Fixed separate vacancy quotas deliberately avoid
asserting substitutability for a firm's complete technology
\cite{kelsocrawford1982}. No global stable or efficient allocation is claimed.

For quote adjustment, $A_{fk}$ counts unique qualified first-choice
applicants to class $k$, not repeat proposals during matching, and
$V_{fk}$ is its initially announced vacancy count. When $V_{fk}>0$,
\begin{equation}
 w^{\rm offer}_{fk,t+1}=\min\left\{w^{\max}_{fk,t},
 w^{\rm offer}_{fk,t}
 \exp\left(\eta_w\frac{V_{fk}-A_{fk}}{V_{fk}+A_{fk}+1}\right)
 \right\}.
 \label{eq:wage}
\end{equation}
Use $\eta_w=0.02$. If $w^{\max}\leq0$, withdraw the vacancy. With no
vacancy, leave the offer quote unchanged. Existing 31-day contracts keep
their agreed wage until renewal. A firm unable to reserve a contracted
day's payroll suspends that shift before work; the failed promise updates
its observed payroll reliability and permits immediate search. There is
no unpaid forced labor or hidden negative cash balance.

\subsection{Homes and leases}
Households rank known affordable vacant dwellings using their complete
plan, including household size, condition, service access, travel, and
Eq.~\eqref{eq:tastes}. A move incurs 10 one-time utility units, two adult
moving hours when available, and actual transport charges. Housing
deposits are zero in the reference model; a later refundable deposit
must be treated as an asset and an escrow liability.

Use household-proposing deferred acceptance, with landlords ranking
applicants by ability to fund the next payment and then stable seeded
priority. Each dwelling is assigned to at most one household and obeys
its occupancy limit. The quote is daily rent under a 30-day lease,
reserved each morning. If payment cannot be funded, the household must
select its available fallback accommodation before using the dwelling;
tenancy protection is a policy extension requiring an identified payer.

For each local dwelling class with $V$ available units and $A$ funded
first-choice applications, update the new-lease quote by
\begin{equation}
 R_{t+1}=R_t\exp\left(\eta_R\frac{A-V}{A+V+1}\right),
 \qquad \eta_R=0.01.
\end{equation}
Occupied units keep their contracted rent. Distinct qualities are
separate classes; a higher quote does not improve a building's physical
attributes.

\subsection{Skills and learning}
Let $s_{ik}\in[0,1]$ describe general, craft, technical, clinical, and
teaching skills. A task with required skill $\bar s_k>0$ receives effective
hours
\begin{equation}
 H_{ik}=h_{ik}(0.6+0.6s_{ik}),
 \label{eq:productivity}
\end{equation}
after enforcing its qualification threshold. For multiple essential
skills use the minimum productivity multiplier. General tasks use general skill without a credential barrier.
The threshold controls eligibility only: skills 0.5, 0.8, and 1 yield
0.90, 1.08, and 1.20 effective hours per actual hour, respectively.

Course progress is attendance times delivered teacher quality times
aptitude. Aptitude satisfies $a_i\in[0.85,1.15]$. For required effective
days $T_e$, progress $P_e$ changes the targeted skill by
\begin{equation}
 s_{ik}=s^{\rm start}_{ik}
 +\min\{P_e/T_e,1\}
 [s^{\rm target}_{ek}-s^{\rm start}_{ik}]_+.
\end{equation}
Completion and assessment confer the credential. Practice outside
courses adds $0.0005(h/8)(1-s)$ per day; an unused occupational skill
loses $(\log 2)/(10\cdot365)s$ daily (approximately a ten-year half-life). General literacy does not decay in the baseline.
Apply learning and decay once per day, then clip to $[0,1]$.

Schools reserve teachers, inputs, and classrooms before offering places.
Applicants rank known programs by Eq.~\eqref{eq:hhplan}; funded
first-choice applications and available new places adjust new tuition
quotes by the rent-style equation with $\eta_T=0.01$. Student-proposing
matching uses eligibility and stable seeded priority, with continuing
students occupying their already contracted places. Publicly financed
tuition is paid from the government account, not omitted from accounting.

An agent estimates employment after training from observed eligible
vacancies and outcomes of known graduates. Let $\hat\pi_e$ be a
Beta$(1,1)$ success/failure estimate for employment within 30 days of
completion, updated only from observed outcomes. Use the median known
eligible wage, or the current fallback wage if no eligible quote is
known; expected income is the $\hat\pi_e$ mixture of qualified and
fallback income. Completion probability is another Beta$(1,1)$ estimate.
These private forecasts can be wrong and can produce education cycles.
No credential carries an assigned monetary premium.

\section{Accounts, resources, and institutional closure}
All households, firms, public agencies, and escrow accounts have
nonnegative cash. With no borrowing or money issue, a transfer of $v$
from $a$ to $b$ changes their balances by $-v$ and $+v$. Taxes, rents,
wages, fees, dividends, and relief are such transfers. Ownership shares
$\theta_{hf}\geq0$ sum to one for each firm. A monthly dividend is
$\theta_{hf}$ times half of cash above a funded seven-day operating
reserve, provided no obligation is overdue. Retained cash stays at the
firm. Liquidation transfers assets; it does not erase money.

Government policies are an input schedule or a declared local rule.
Taxes are collected only from their specified realized base; subsidies
and benefits cannot exceed available treasury cash. If a payment cannot
be funded, the policy is rationed by a published priority rule. The
reference experiment can set every tax and transfer to zero. A change
in public policy must not silently create goods, workers, or currency.

For each storable good let $I_g$ count physical lots, including transit
and household stocks but excluding installed assets and consumed batch
inputs. Over the half-open interval between successive start-of-day states,
\begin{align}
 I_{g,t+1}={}&I_{g,t}+B^{\rm complete}_{g,t}
 +B^{\rm well}_{g,t}\nonumber\\
 &-A^{\rm batch}_{g,t}-A^{\rm venue}_{g,t}
 -A^{\rm power}_{g,t}\nonumber\\
 &-C^{\rm food/home}_{g,t}-A^{\rm build}_{g,t}
 -A^{\rm repair}_{g,t}\nonumber\\
 &-A^{\rm install}_{g,t}-W^{\rm retire}_{g,t}
 -W^{\rm expiry}_{g,t}.
 \label{eq:material}
\end{align}
Venue inputs include bathhouse water/soap, clinic medicines, school inputs
and all delivered-service coefficients. Home use includes food, drinking
and washing water, soap and cooking fuel/water. Construction, repair and
installation have separate event codes, each debited once. Well output
is explicit. Power fuel is debited as delivered flow is produced.
Transfers between owners cancel. Durable use changes a separate condition
ledger; only retirement removes the physical item from $I_g$. Installed
asset counts, construction progress, condition and retired waste have
separate ledgers; condition-weighted value never replaces a lot count.
Recipe tool wear is classified as batch input, not retirement or repair.
Electricity satisfies delivered generation equals delivered use (zero
storage and zero network loss in this baseline) and is excluded from this
inventory equation.

Natural-resource regeneration is a scenario parameter. A renewable fish
stock may use $Z_{t+1}=Z_t+rZ_t(1-Z_t/K)-q_t$, with harvest restricted
to available stock. Forest cohorts mature after the specified delay.
These physical constraints do not assign a price to nature. Land can be
privately leased or publicly allocated, but ownership and use must be
specified. The closed reference economy has no unlimited export buyer.

\section{Analytical properties and their limits}
\begin{theorem}[Accounting and physical feasibility]
Suppose initial cash and inventories are nonnegative, reservations are
atomic, allocations satisfy the stated capacity and time constraints,
and every execution is bounded by its reservation. Then the reference
algorithm preserves nonnegative cash and inventories, prevents duplicate
use of committed goods and time, and releases material output only after
its production requirements have been completed. In the closed
no-credit economy, total money including escrow is conserved.
\end{theorem}
\begin{proof}
Proceed by induction through the event order. A reservation reduces
spendable resources without changing their total ownership. Since all
orders of an account share the same ledger, their costs cannot exceed
its free balance. Equations~\eqref{eq:buyfill}--\eqref{eq:sellfill}
execute no more than each reserved order, and aggregate purchases and
sales both equal $Q_m$. Each payment is an equal debit and credit,
including transfers into and out of escrow. Consumption, expiry,
maintenance, and input commitment are bounded by owned available stock.
Slots and staff schedules are reserved before execution. Equation
\eqref{eq:progress} cannot exceed remaining progress, labor, electricity,
or the process rate, and completion is required before output release.
Refunding unused reservations reverses no completed transaction. Thus
each event preserves the induction hypotheses, and no event changes
the sum of cash balances. Monetary conservation does not imply
conservation of utility, employment, or the quantity of each good.
\end{proof}

\begin{lemma}[Finite-horizon price bounds]
For a market with constant $\eta>0$ and initial $p_0>0$,
\begin{equation}
 p_0e^{-\eta t}\leq p_t\leq p_0e^{\eta t}.
\end{equation}
\end{lemma}
\begin{proof}
The normalized imbalance has absolute value at most one. Sum the
logarithmic increments in Eq.~\eqref{eq:price}.
\end{proof}
These bounds prevent a finite-step zero or infinity in exact arithmetic.
They give no uniform infinite-horizon bound. Numerical implementations
should store log quotes; a numerical underflow guard is not an economic
price floor.

\begin{lemma}[Information locality]
If initial beliefs are local and independent matching randomness is
recorded, every action defined above is measurable with respect to the
agent's information and its received local messages. The algorithm
requires no counterfactual global production allocation.
\end{lemma}
\begin{proof}
Household plans use their own states, known candidate sets, and local
beliefs. Firms use their own recipes, finance, inventories, workers, and
market announcements. Clearing uses submitted orders only. All
subsequent messages are functions of these quantities and recorded
randomness. Induction on event time establishes the claim.
\end{proof}

\begin{proposition}[Nominal scale equivariance]
Multiply initial cash, all quotes and contracts, monetary histories, and
monetary policy thresholds by the same $c>0$. Preserve physical states,
dimensionless parameters, and random draws. In exact arithmetic, the
algorithm produces the same real allocations, with every monetary
quantity multiplied by $c$.
\end{proposition}
\begin{proof}
Budget and escrow inequalities scale by $c$, so their feasible real
bundles are unchanged. The essential-basket cost and the finite-difference
step in Eq.~\eqref{eq:lambda} scale by $c$, while $\lambda$ scales by
$1/c$; the terminal utility is therefore unchanged. Firm objectives and
maximum profitable wages scale by $c$ without changing their ranking.
Tie rules use scale-invariant comparisons. Funded quantities, matchings,
and physical updates are unchanged. Normalized imbalances are unchanged,
so multiplicative price, wage, rent, and tuition updates preserve the
scaling. Induct over days.
\end{proof}
Changing the distribution of wealth is not this transformation. Neither
path independence nor a unique relative-price vector follows.

\subsection{Local price stability}
Freeze other states and suppose a differentiable price-only response
$e(\ell)$ exists for log quotes $\ell=\log p$. This assumption fails at
many discrete switches and must be checked in the particular benchmark.
At $e(\ell^*)=0$, put $J=De(\ell^*)$. For heterogeneous gains the
linearization is $I+\operatorname{diag}(\eta_m)J$, and its spectral radius
must be less than one. With a common step $\eta$,
\begin{equation}
 \delta\ell_{t+1}=(I+\eta J)\delta\ell_t
 +o(\norm{\delta\ell_t}).
 \label{eq:linear}
\end{equation}
\begin{proposition}[Conditional local stability]
If $\rho(I+\eta J)<1$, the fixed point is locally asymptotically
stable. If an eigenvalue has modulus greater than one, it is unstable.
Eigenvalues on the unit circle leave the linearization test inconclusive.
For $\lambda\in\operatorname{spec}(J)$ with
$\operatorname{Re}\lambda<0$, its linear mode is stable precisely when
\begin{equation}
 0<\eta<\frac{-2\operatorname{Re}\lambda}{|\lambda|^2}.
\end{equation}
\end{proposition}
\begin{proof}
The first assertions are the discrete-time linearization criterion.
The last follows by expanding $|1+\eta\lambda|^2<1$.
\end{proof}
A positive-real-part eigenvalue cannot be stabilized merely by reducing
a positive step. For the actual economy one must augment the Jacobian
with inventories, finance, contracts, expectations, and pipelines. The
price-only proposition is not a proof for that augmented system.

\subsection{A separate convex coordination benchmark}
To isolate a provable price mechanism, consider compact convex sets
$Y_a$, concave continuous payoffs $f_a$, and affine resource use $r_a$.
Let
\begin{equation}
 \max_{y_a\in Y_a}\sum_a f_a(y_a)
 \quad\text{subject to}\quad\sum_a r_a(y_a)\leq b.
 \label{eq:benchmark}
\end{equation}
Assume Slater feasibility; compactness makes the primal optimum finite,
and Slater's condition gives dual attainment. Interpret $f_a$
as quasilinear net benefit in a common numeraire with nonbinding finance;
it is not an arbitrary sum of the original households' cardinal utilities.
The dual is
\begin{equation}
 \Phi(p)=p^\top b+
 \sum_a\max_{y_a\in Y_a}\{f_a(y_a)-p^\top r_a(y_a)\},
 \quad p\geq0.
\end{equation}
Each maximization is private. Resource markets need only summed resource
use, consistent with dual decomposition \cite{boyd2004,agrawal2022}.

\begin{theorem}[Restricted dual-price convergence]
Under the preceding assumptions and
positive steps with $\sum_t\alpha_t=\infty$ and
$\sum_t\alpha_t^2<\infty$, the update
\begin{equation}
 p_{t+1}=\left[p_t+\alpha_t
 \left(\sum_a r_a(y_{a,t})-b\right)\right]_+
 \label{eq:dualupdate}
\end{equation}
with exact private maximizers converges to a dual minimizer in finite
dimension.
\end{theorem}
\begin{proof}
A subgradient of $\Phi$ is $g_t=b-\sum_a r_a(y_{a,t})$.
Compactness of the finite product of $Y_a$ and affine resource maps gives
$\norm{g_t}\leq G<\infty$.
Nonexpansiveness of projection and the subgradient inequality give, for
any dual optimum $p^*$,
\begin{align}
 \norm{p_{t+1}-p^*}^2\leq{}&\norm{p_t-p^*}^2
 -2\alpha_t[\Phi(p_t)-\Phi(p^*)]\nonumber\\
 &+\alpha_t^2G^2.
\end{align}
Summability of the last term makes distances to the optimum set
convergent and the sequence bounded. Since the steps have divergent
sum, the nonnegative objective gap has zero liminf. A bounded
subsequence attaining that liminf has an optimal cluster point by
continuity. Distance convergence to this cluster point then implies
convergence of the whole sequence.
\end{proof}
This is an additive projected dual update for a different, restricted
problem. It does not prove convergence of the normalized multiplicative
cash-constrained algorithm, convergence of every primal allocation, or
efficiency of the game's distribution of income.

\section{Implementation defaults and validation protocol}
Physical catalogs and preferences are scenario inputs, fixed before an
experiment. Use at least 35 days of initial basic food, working assets,
essential skilled workers, and viable water and fuel access. Test whether
crop plots and initial production capacity can physically support the
population before attributing starvation to a market rule. Fixed initial
assets are necessary to bootstrap circular tool--steel--machinery chains.

Record the simulator version, scenario checksum, random seed, price
initialization, wealth distribution, and every parameter. Stagger review
dates by a seeded uniform offset, rather than letting every household or
firm adjust simultaneously. Default gains and review periods are listed
in Appendix~\ref{app:parameters}; they are experimental values, not
established stabilizing constants.

The required diagnostic suite is: (i) exact cash and reservation checks;
(ii) physical flow and age bookkeeping; (iii) nominal rescaling;
(iv) consumer utility and substitution checks at fixed quotes;
(v) local stability analysis where differentiability applies; and
(vi) repeated-seed shock experiments with harvest, fuel, workforce,
transport, and schooling disruptions. Report unmet needs, unemployment,
vacancies, service queues, stranded complementary inputs, bankruptcies,
price dispersion, and recovery times. Failed economies are outcomes to
report, not runs to discard.

Compare alternative market rules while holding physical endowments,
information, and institutions fixed. A centralized allocation solver may
serve as an external benchmark, but its solution must not be passed to
agents. Distinguish financial insolvency, physical impossibility, search
failure, and coordination failure. A high price alone identifies none
of these uniquely.

The reported checks comprise food-processing nutrient inequalities and
the static feasibility calculation in Appendix~\ref{app:feasibility}. No full-economy simulation results, estimated
parameters, convergence rates, or empirical fit are reported here.

The reference daily discount gives $\beta^{365}\simeq0.9642$ and
$\beta^{720}\simeq0.9305$, an annual effective rate of about 3.72\%.
Goods can adjust every day, new wages every seven days, and new rents
and tuition every thirty days; maximal average log-adjustment speeds
are therefore 0.02, 0.02/7 and 0.01/30 per day. This is an explicit
contractual rigidity hypothesis. It can generate real-wage swings.
Sensitivity experiments must also use common per-day speeds, e.g.
review gains $7\eta$ and $30\eta$, and vary the common $\eta$; the
reference asymmetry is not a stability calibration.

\section{Conclusion}
The model specifies a local monetary mechanism through which household
needs generate retail orders, firms transmit derived orders upstream,
and wages, rents, and education respond to locally observed opportunities.
Physical production and service capacity prevent prices from substituting
for missing resources. Explicit reservation and settlement establish
feasibility independently of whether the economy is prosperous or stable.
The main unresolved scientific question is the stability and efficiency
of the coupled adaptive system under complementary technologies,
financing constraints, and delayed skill supply. The proposed catalog and
algorithm make that question reproducible without presuming its answer.



% Wide multipage reference tables appear in a full-width appendix.
\onecolumn
\appendix
\section{Goods, units, and storage}\label{app:goods}
The catalog comprises exactly 34 commodity classes. One water unit is
100 liters and one electricity unit is 10 kWh. Other quantities are
commodity-specific lots. In particular, a grain lot supplies one
adult-day of the energy index before losses; it is not interchangeable
with a kilogram of every other good. Equipment and garments are durable
lots; fractional orders can represent pooled procurement, while household
allocation and physical objects can be made indivisible in a declared
extension. No row has a base price.

Goods are tracked by lot age and condition. Shelf-life expiry is an
explicit removal to waste, not an unexplained second depreciation charge.
Waste remains a recorded physical residual. The reference scenario
provides a finite permitted landfill volume as an endowment; once full,
waste-producing activity requires funded disposal capacity or stops.
No disposal industry or pollution-removal technology is assumed to exist
without an explicit recipe. Electricity cannot be stockpiled.

\begingroup
\footnotesize
\setlength{\tabcolsep}{2pt}
\renewcommand{\arraystretch}{1.15}
\begin{longtable}{@{}>{\raggedright\arraybackslash}p{0.05\textwidth}>{\raggedright\arraybackslash}p{0.18\textwidth}>{\raggedright\arraybackslash}p{0.34\textwidth}>{\raggedright\arraybackslash}p{0.34\textwidth}@{}}
\caption{Launch commodity catalog and storage conventions.}\label{tab:goods}\\
\toprule
\textbf{ID} & \textbf{Commodity} & \textbf{Uses} & \textbf{Storage / durability} \\
\midrule
\endfirsthead
\multicolumn{4}{l}{\textit{Table \thetable{} (continued)}}\\
\toprule
\textbf{ID} & \textbf{Commodity} & \textbf{Uses} & \textbf{Storage / durability} \\
\midrule
\endhead
\midrule
\multicolumn{4}{r}{\textit{Continued on next page}}\\
\endfoot
\bottomrule
\endlastfoot
G01 & Grain & Milling; household staple preparation & 180-day shelf life \\
G02 & Legumes & Household food; preserved food; prepared meals & 180 days \\
G03 & Vegetables and fruit & Household food; diet quality; food processing & 7 days \\
G04 & Fish & Household protein; food processing & 2 days without refrigeration \\
G05 & Timber & Boards; charcoal; paper & Dry storage; no spoilage in the reference model \\
G06 & Textile fibre & Fabric & Dry storage \\
G07 & Oilseed & Edible oil & 180 days \\
G08 & Iron ore & Steel & No spoilage \\
G09 & Limestone & Cement; steel flux & No spoilage \\
G10 & Clay & Bricks; cement & No spoilage \\
G11 & Salt & Preserving food; soap & Dry storage \\
G12 & Medicinal plants & Medicines & 30 days \\
G13 & Flour & Bread; household staple preparation & 90 days \\
G14 & Boards & Buildings; furniture; tools; machinery & Dry storage \\
G15 & Fabric & Garments; furniture; medicines & Dry storage \\
G16 & Charcoal & Process heat; cooking; power & Dry storage \\
G17 & Steel & Machinery; tools; construction; food containers & No spoilage \\
G18 & Cement & Construction and maintenance & 90 days in dry storage \\
G19 & Bricks & Construction and maintenance & No spoilage \\
G20 & Edible oil & Cooking; meals; soap & 90 days \\
G21 & Paper & Books; schools; administration; packaging & Dry storage \\
G22 & Tool kits & Equipment and maintenance inputs & Durability stock, not an edible flow \\
G23 & Machinery modules & Plant capacity and repair & Durability stock \\
G24 & Soap & Household hygiene; clinics; services & 365 days \\
G25 & Bread & Convenient staple food & 3 days \\
G26 & Preserved food & Durable mixed food & 365 days \\
G27 & Prepared meals & Convenient, pleasant mixed food & 1 day \\
G28 & Basic clothing & Coverage and basic comfort & 180 wear-days at full condition \\
G29 & Fine clothing & Coverage, comfort, and style & 240 wear-days \\
G30 & Furniture kits & Household comfort; venue equipment & 1,800 occupied days \\
G31 & Books & Reading; libraries; education & 180 reading sessions before replacement \\
G32 & Medicines & Treatment input; no direct consumption pleasure & 180 days \\
G33 & Clean water & Drinking; hygiene; processing & Stored water lasts 7 days unless treated \\
G34 & Electricity & Machines; lighting; ventilation; services & Instantaneous network service; no storage in the reference model \\
\end{longtable}
\endgroup

\section{Production coefficients and construction}\label{app:recipes}
The following tables provide all commodity-production routes in the
launch catalog. The hours column is total effective labor per batch;
the duration is a minimum process time. Worker and slot limits are
separate constraints. Input quantities are game parameters, not measured
industrial coefficients. An asterisk identifies replacement/tool wear
consumed from inventory; it does not replace the need for installed
equipment. Material output is released at the next boundary after completion.

Any required named supervisor must be a qualified member of the listed
workforce. For a recipe requiring one or two technicians, allocate one
or two equal worker-shares of its labor requirement to technical labor:
one share is total recipe hours divided by the declared staffing reference (the worker cap). This fixed labor mix is
independent of realized employment. The rest is
general labor. A craft-plus-technician machine shop allocates one such
share to technical labor and the rest to craft labor. A medicinal-garden
supervisor requires craft skill at least 0.5. This convention makes the
labor vectors in the main text numerical.

\begingroup
\footnotesize
\setlength{\tabcolsep}{2pt}
\renewcommand{\arraystretch}{1.15}
\begin{longtable}{@{}>{\raggedright\arraybackslash}p{0.145\textwidth}>{\raggedright\arraybackslash}p{0.305\textwidth}>{\raggedright\arraybackslash}p{0.105\textwidth}>{\raggedright\arraybackslash}p{0.04\textwidth}>{\raggedright\arraybackslash}p{0.07\textwidth}>{\raggedright\arraybackslash}p{0.085\textwidth}>{\raggedright\arraybackslash}p{0.14\textwidth}@{}}
\caption{Extraction and agriculture.}\label{tab:recipes0}\\
\toprule
\textbf{Facility} & \textbf{Inputs} & \textbf{Output} & \textbf{h} & \textbf{Lag} & \textbf{Workers / slots} & \textbf{Skills} \\
\midrule
\endfirsthead
\multicolumn{7}{l}{\textit{Table \thetable{} (continued)}}\\
\toprule
\textbf{Facility} & \textbf{Inputs} & \textbf{Output} & \textbf{h} & \textbf{Lag} & \textbf{Workers / slots} & \textbf{Skills} \\
\midrule
\endhead
\midrule
\multicolumn{7}{r}{\textit{Continued on next page}}\\
\endfoot
\bottomrule
\endlastfoot
Grain farm & 2 water + 0.01 tools* & 20 grain & 40 & 30 days & 2 / 12 plots & General \\
Legume farm & 2 water + 0.01 tools* & 10 legumes & 32 & 30 days & 2 / 12 plots & General \\
Vegetable farm & 4 water + 0.01 tools* & 20 vegetables & 32 & 10 days & 2 / 4 plots & General \\
Fibre farm & 3 water + 0.01 tools* & 10 fibre & 32 & 30 days & 2 / 12 plots & General \\
Oilseed farm & 2 water + 0.01 tools* & 10 oilseed & 32 & 30 days & 2 / 12 plots & General \\
Medicinal garden & 2 water + 0.01 tools* & 5 medicinal plants & 24 & 20 days & 4 / 6 plots & General; one trained supervisor \\
Fishing dock & 0.02 tools* & 10 fish & 32 & 1 day & 8 / 2 boats & General \\
Forestry camp & 0.03 tools* & 10 timber & 24 & 1 day & 8 / 3 cutting slots & General \\
Iron mine & 0.05 tools* + 0.1 electricity & 10 ore & 40 & 1 day & 12 / 3 slots & General; one technician \\
Limestone quarry & 0.03 tools* & 10 limestone & 24 & 1 day & 8 / 3 slots & General \\
Clay pit & 0.02 tools* & 10 clay & 16 & 1 day & 6 / 3 slots & General \\
Saltworks & Seawater access + 0.01 tools* & 5 salt & 16 & 3 days & 6 / 9 evaporation beds & General \\
\end{longtable}
\endgroup

\begingroup
\footnotesize
\setlength{\tabcolsep}{2pt}
\renewcommand{\arraystretch}{1.15}
\begin{longtable}{@{}>{\raggedright\arraybackslash}p{0.145\textwidth}>{\raggedright\arraybackslash}p{0.305\textwidth}>{\raggedright\arraybackslash}p{0.105\textwidth}>{\raggedright\arraybackslash}p{0.04\textwidth}>{\raggedright\arraybackslash}p{0.07\textwidth}>{\raggedright\arraybackslash}p{0.085\textwidth}>{\raggedright\arraybackslash}p{0.14\textwidth}@{}}
\caption{Intermediate industry and utilities.}\label{tab:recipes1}\\
\toprule
\textbf{Facility} & \textbf{Inputs} & \textbf{Output} & \textbf{h} & \textbf{Lag} & \textbf{Workers / slots} & \textbf{Skills} \\
\midrule
\endfirsthead
\multicolumn{7}{l}{\textit{Table \thetable{} (continued)}}\\
\toprule
\textbf{Facility} & \textbf{Inputs} & \textbf{Output} & \textbf{h} & \textbf{Lag} & \textbf{Workers / slots} & \textbf{Skills} \\
\midrule
\endhead
\midrule
\multicolumn{7}{r}{\textit{Continued on next page}}\\
\endfoot
\bottomrule
\endlastfoot
Mill & 10 grain + 0.1 electricity & 9 flour & 8 & 1 day & 4 / 4 & General \\
Sawmill & 10 timber + 0.1 electricity + 0.02 tools* & 8 boards & 12 & 1 day & 6 / 4 & General \\
Textile mill & 10 fibre + 1 water + 0.2 electricity & 8 fabric & 24 & 1 day & 8 / 3 & General; one technician \\
Charcoal kiln & 10 timber & 4 charcoal & 8 & 2 days & 4 / 8 & General \\
Steelworks & 10 ore + 5 charcoal + 1 limestone + 1 water + 0.2 electricity & 5 steel & 40 & 2 days & 12 / 5 & General; two technicians \\
Cement works & 8 limestone + 2 clay + 2 charcoal + 0.2 electricity & 8 cement & 24 & 2 days & 8 / 5 & General; one technician \\
Brickworks & 10 clay + 1 water + 2 charcoal & 8 bricks & 16 & 2 days & 8 / 6 & General \\
Oil press & 10 oilseed + 0.1 electricity & 2.5 oil & 12 & 1 day & 4 / 2 & General \\
Paper mill & 4 timber + 4 water + 1 charcoal + 0.2 electricity & 8 paper & 24 & 2 days & 8 / 5 & General; one technician \\
Tool workshop & 2 steel + 1 boards + 0.1 electricity & 1 tools & 16 & 1 day & 6 / 3 & Craft \\
Machine shop & 5 steel + 1 boards + 0.1 tools* + 0.5 electricity & 1 machinery & 40 & 3 days & 8 / 4 & Craft plus technician \\
Soap works & 1 oil + 0.2 salt + 0.5 water + 0.2 charcoal & 3 soap & 8 & 1 day & 4 / 3 & Craft \\
Water plant & Freshwater access + 0.1 electricity & 10 water & 2 & 1 day & 4 / 12 & One technician \\
Power station & 2 charcoal + 0.2 water & 10 electricity & 4 & Flow & 6 / 12 flow units & Two technicians \\
\end{longtable}
\endgroup

\begingroup
\footnotesize
\setlength{\tabcolsep}{2pt}
\renewcommand{\arraystretch}{1.15}
\begin{longtable}{@{}>{\raggedright\arraybackslash}p{0.145\textwidth}>{\raggedright\arraybackslash}p{0.305\textwidth}>{\raggedright\arraybackslash}p{0.105\textwidth}>{\raggedright\arraybackslash}p{0.04\textwidth}>{\raggedright\arraybackslash}p{0.07\textwidth}>{\raggedright\arraybackslash}p{0.085\textwidth}>{\raggedright\arraybackslash}p{0.14\textwidth}@{}}
\caption{Food processing and household goods.}\label{tab:recipes2}\\
\toprule
\textbf{Facility} & \textbf{Inputs} & \textbf{Output} & \textbf{h} & \textbf{Lag} & \textbf{Workers / slots} & \textbf{Skills} \\
\midrule
\endfirsthead
\multicolumn{7}{l}{\textit{Table \thetable{} (continued)}}\\
\toprule
\textbf{Facility} & \textbf{Inputs} & \textbf{Output} & \textbf{h} & \textbf{Lag} & \textbf{Workers / slots} & \textbf{Skills} \\
\midrule
\endhead
\midrule
\multicolumn{7}{r}{\textit{Continued on next page}}\\
\endfoot
\bottomrule
\endlastfoot
Bakery & 10 flour + 3 water + 1 charcoal & 14 bread & 12 & 1 day & 6 / 4 & Craft \\
Food preservatory & 6 fish + 4 legumes + 3 vegetables + 0.5 oil + 0.3 salt + 0.2 steel + 0.2 paper + 1 charcoal + 1 water & 10 preserved food & 24 & 1 day & 8 / 3 & Craft \\
Kitchen / restaurant & 4 bread + 3 legumes + 4 vegetables + 2 fish + 0.5 oil + 1 water + 0.5 charcoal & 10 prepared meals & 16 & 1 day & 8 / 4 daily batches & Craft \\
Clothing workshop, basic & 2 fabric & 1 basic clothing & 8 & 1 day & 8 / 8 & Craft \\
Clothing workshop, fine & 3 fabric + 0.02 tools* & 1 fine clothing & 16 & 2 days & 8 / 8 & Advanced craft \\
Furniture workshop & 3 boards + 0.5 fabric + 0.1 steel + 0.02 tools* & 1 furniture & 16 & 2 days & 8 / 8 & Craft \\
Printer & 1 paper + 0.05 oil + 0.05 electricity & 1 books & 2 & 1 day & 4 / 12 & Literate craft \\
Pharmacy works & 1 medicinal plants + 0.2 fabric + 0.1 paper + 0.2 water + 0.05 electricity & 5 medicines & 8 & 2 days & 4 / 6 & Clinical technician \\
\end{longtable}
\endgroup

Freshwater, seawater, agricultural plots, fish biomass, forest cohorts,
and mineral deposits are physical site resources. Crop yields multiply
by fertility in $[0,1]$. Water is an upfront input; a batch lacking its
full water requirement cannot start. There is no separate water-completion factor. Crop-farm worker caps are two: at unit productivity a grain farm's
12 plots require $12(40)/30=16$ hours daily, exactly two eight-hour
workers. Lower skill or absence makes labor binding; faster workers
remain limited by plots. Farm labor is entirely general, so no
supervisor-share rule is applied to these farms. Farm plots are occupied for
the entire growth cycle. Forestry consumes standing timber and regrows
by cohorts after 1,800 days. Fish use logistic regeneration with scenario
parameters $r$ and $K$; mineral extraction reduces finite deposits.

A non-electric emergency well yields one clean-water unit per four
actual labor-hours at a suitable source. The soap recipe aggregates
minor mineral reagents into the salt input; it is not a chemical
manufacturing instruction. All unlisted residues are waste, not free
additional output.

For construction, letters A--E below denote bills of materials, not
tradable commodities or currency-valued building costs. Land and access
are additional constraints. A construction worker uses an installed
tool allocation of one kit per four workers; wear consumes 0.01 tool
kits per 80 actual construction hours.

\begingroup
\footnotesize
\setlength{\tabcolsep}{2pt}
\renewcommand{\arraystretch}{1.15}
\begin{longtable}{@{}>{\raggedright\arraybackslash}p{0.34\textwidth}>{\raggedright\arraybackslash}p{0.08\textwidth}>{\raggedright\arraybackslash}p{0.08\textwidth}>{\raggedright\arraybackslash}p{0.08\textwidth}>{\raggedright\arraybackslash}p{0.08\textwidth}>{\raggedright\arraybackslash}p{0.1\textwidth}>{\raggedright\arraybackslash}p{0.08\textwidth}@{}}
\caption{Structural construction packages.}\label{tab:build}\\
\toprule
\textbf{Package} & \textbf{Boards} & \textbf{Bricks} & \textbf{Cement} & \textbf{Steel} & \textbf{h} & \textbf{Days} \\
\midrule
\endfirsthead
\multicolumn{7}{l}{\textit{Table \thetable{} (continued)}}\\
\toprule
\textbf{Package} & \textbf{Boards} & \textbf{Bricks} & \textbf{Cement} & \textbf{Steel} & \textbf{h} & \textbf{Days} \\
\midrule
\endhead
\midrule
\multicolumn{7}{r}{\textit{Continued on next page}}\\
\endfoot
\bottomrule
\endlastfoot
A: light structure & 12 & 4 & 0 & 0 & 80 & 4 \\
B: small masonry structure & 12 & 20 & 8 & 2 & 160 & 8 \\
C: workshop / service building & 20 & 40 & 16 & 5 & 320 & 12 \\
D: industrial hall & 30 & 80 & 40 & 20 & 640 & 20 \\
E: apartment block / large institution & 40 & 160 & 80 & 40 & 1,280 & 30 \\
\end{longtable}
\endgroup

Facility installation packages are: farms A plus 0.2 machinery modules;
medicinal garden A plus 0.1; fishing dock B plus 1;
forestry, quarries, clay pits, and saltworks A plus 0.5;
mine D plus 3; mill, sawmill, and oil press C plus 2;
textile and paper mills D plus 4; kiln and brickworks C plus 1;
steel and cement works D plus 6; tools, garments, furniture, soap,
bakery, and printer C plus 1; machine shop D plus 4;
preservatory D plus 3; pharmacy C plus 2; water and power plants D plus 5.
Each industrial establishment also installs at least one tool kit.

Installed machinery requires replacements equal to 0.0005 of installed
modules per operating day. Each building, occupied or vacant, requires materials
and repair labor equal to 0.0001 of its original construction bill daily.
Accumulate fractional repair credits rather than transporting infinitesimal
pieces. Let $r_t$ be delivered repair as a fraction of the full original
bill. Set daily building depreciation $d=10^{-4}$, including vacancy;
machinery depreciation is $5\cdot10^{-4}$ on operating days and zero
when mothballed. Update condition as
$c_{t+1}=\min\{1,\max\{0,c_t-d\}+r_t\}$; all required repair components
must be supplied proportionally. Machinery uses the analogous rule with
its own replacement rate. Usable capacity and shelter multiply by
condition. Replacements restore condition, not additional installed
capacity. Tools explicitly consumed by a recipe are not charged a second
maintenance debit for that use.

\section{Consumption, utility, and health parameters}\label{app:utility}
Food attributes are defined per consumed lot. The adult targets are
energy $N=1$, protein $P=0.5$, and micronutrients $M=0.5$; $Q$ is sensory
quality. These are simulation indexes, not dietary recommendations.

\begingroup
\footnotesize
\setlength{\tabcolsep}{2pt}
\renewcommand{\arraystretch}{1.15}
\begin{longtable}{@{}>{\raggedright\arraybackslash}p{0.28\textwidth}>{\raggedright\arraybackslash}p{0.1\textwidth}>{\raggedright\arraybackslash}p{0.1\textwidth}>{\raggedright\arraybackslash}p{0.1\textwidth}>{\raggedright\arraybackslash}p{0.13\textwidth}>{\raggedright\arraybackslash}p{0.2\textwidth}@{}}
\caption{Food-service coefficients.}\label{tab:food}\\
\toprule
\textbf{Food} & \textbf{Energy N} & \textbf{Protein P} & \textbf{Micro. M} & \textbf{Enjoyment Q} & \textbf{Family} \\
\midrule
\endfirsthead
\multicolumn{6}{l}{\textit{Table \thetable{} (continued)}}\\
\toprule
\textbf{Food} & \textbf{Energy N} & \textbf{Protein P} & \textbf{Micro. M} & \textbf{Enjoyment Q} & \textbf{Family} \\
\midrule
\endhead
\midrule
\multicolumn{6}{r}{\textit{Continued on next page}}\\
\endfoot
\bottomrule
\endlastfoot
Grain & 1.00 & 0.15 & 0.05 & 0.05 & Staple \\
Legumes & 0.80 & 0.70 & 0.20 & 0.15 & Legume \\
Vegetables / fruit & 0.25 & 0.10 & 1.00 & 0.25 & Produce \\
Fish & 1.00 & 1.00 & 0.30 & 0.30 & Fish \\
Flour & 1.10 & 0.15 & 0.02 & 0.05 & Staple \\
Edible oil & 3.00 & 0 & 0 & 0.10 & Oil \\
Bread & 0.75 & 0.10 & 0.01 & 0.35 & Staple \\
Preserved food & 1.00 & 0.80 & 0.40 & 0.35 & Mixed \\
Prepared meals & 0.95 & 0.45 & 0.45 & 1.00 & Mixed \\
\end{longtable}
\endgroup

Grain, flour, legumes, and raw fish require preparation. Preparing one
energy unit uses 0.4 adult-hours, 0.03 charcoal, and 0.01 water, in
addition to the food itself. Serving bread or preserved food uses 0.05
hours per energy unit; purchased prepared meals require no further
preparation. Eating time is already in personal care. Oil is an ingredient
limited to 0.10 lots per adult-day and must accompany other food.

The four food-processing recipes satisfy the componentwise inequalities
\[A^{\rm nutrient}b_r\leq A^{\rm nutrient}a_r.\]
The output/input energy ratios for milling, baking, preserving, and
meal preparation are respectively $0.990$, $0.955$, $0.873$, and $0.960$.
Enjoyment is not a conserved nutrient and may increase through cooking.

One wearable clothing set covers one adult irrespective of remaining
condition. One wear-day removes $1/180$ condition from a basic set or
$1/240$ from a fine set. Only one set may be assigned; fine style equals
its remaining condition.
Furniture loses $1/1800$ condition per
occupied day. A book copy lasts 180 reading sessions. Title-specific
novelty multiplies reading service by $1,1/2,1/4$, then $1/8$ on later
readings. Library lending reserves the physical copy.

\begingroup
\footnotesize
\setlength{\tabcolsep}{2pt}
\renewcommand{\arraystretch}{1.15}
\begin{longtable}{@{}>{\raggedright\arraybackslash}p{0.22\textwidth}>{\raggedright\arraybackslash}p{0.63\textwidth}>{\raggedright\arraybackslash}p{0.07\textwidth}@{}}
\caption{Baseline weights in Eq.~\eqref{eq:utility}.}\label{tab:weights}\\
\toprule
\textbf{Service} & \textbf{Normalization / interpretation} & \textbf{Weight} \\
\midrule
\endfirsthead
\multicolumn{3}{l}{\textit{Table \thetable{} (continued)}}\\
\toprule
\textbf{Service} & \textbf{Normalization / interpretation} & \textbf{Weight} \\
\midrule
\endhead
\midrule
\multicolumn{3}{r}{\textit{Continued on next page}}\\
\endfoot
\bottomrule
\endlastfoot
Safe drinking water & $d_i/0.03$ & 120 \\
Energy from food & $N_i$ & 100 \\
Weatherproof shelter & occupied home's shelter index & 50 \\
Clothing coverage & one assigned wearable set & 28 \\
Protein & $P_i/0.5$ & 18 \\
Micronutrients & $M_i/0.5$ & 18 \\
Hygiene & home hygiene + purchased bathhouse hygiene & 14 \\
Thermal comfort & climate-adjusted comfort index & 14 \\
Space & $A_h/(20\,n_h^{\rm eq}\,\mathrm{m}^2)$ & 8 \\
Privacy & household privacy index & 6 \\
Furnishing & effective furniture kits / adult equivalents & 5 \\
Clean environment & 1 - residential pollution exposure & 12 \\
Personal security & 1 - local crime exposure & 10 \\
Social connection & social stock & 10 \\
Recreation & recreation stock & 8 \\
Reading / cultural engagement & reading stock & 5 \\
Religious practice & religious stock & 0 to 8 \\
Food enjoyment & $\sum_g\zeta_{ig}q_gx_{ig}$ & 6 \\
Clothing style & condition of the assigned fine set & 4 \\
\end{longtable}
\endgroup

Multiply non-survival weights independently by a seeded draw from
$[0.85,1.15]$; water and food-energy weights instead use $[0.95,1.05]$.
Half the initial population has religious weight zero; the remainder
draw uniformly from $[2,8]$ without a further multiplier. These are synthetic heterogeneity assumptions.
No food-variety entropy bonus is included.

For a minimal fictional health module, with $d_t$ drinking water in
water units, use
\begin{equation}
 h_{t+1}=\operatorname{clip}_{[0,100]}\!\left[
 h_t+0.2-4(1-N_t)_+-8(1-d_t/0.03)_+
 -I_t^{\rm loss}-J_t^{\rm loss}+T_t^{\rm gain}\right].
\end{equation}
A minor illness lasts seven symptomatic days and causes one health-point
loss daily. Only a currently uninfected person draws a new illness, with
probability $0.002+0.01(1-\min\{z_{\rm hygiene},1\})+
0.005\,\text{pollution exposure}$. One successful clinic visit reduces
remaining symptomatic time by $3q$ days with probability 0.8; fractional
days reduce the last day's loss proportionally. Same-day repeat treatment
has no further effect. The clinic cannot improve health above 100 by
selling visits to a healthy person.

Per work-hour, a lost-time injury occurs with probability
$2\times10^{-6}+2\times10^{-5}r$. Draw one uniform integer absence
duration from 3 to 10 days and an immediate 10-point health loss; at most
one new injury is drawn per person-day. Injured persons do not work during
their absence. This hazard module is illustrative, not occupational
epidemiology. Severe illness and hospital treatment are disabled in the
reference stochastic experiment until a treatment transition is supplied;
the hospital remains a catalogued service technology for that extension.

Residential pollution is produced by operating batches: 0.05 load for
farms and fishing, 0.2 for extraction and light industry, and 1.0 for
charcoal, steel, cement, and power. Allocate each load among cells using
normalized weights $\exp(-d/1\,\mathrm{km})$, with distance measured
from the emitting establishment. Local stock retains 90\% daily and
exposure is $1-\exp(-\mathrm{stock}/10)$. Services contribute 0.02 per
operating establishment-day. These are synthetic environmental units.
Crime exposure is a fixed scenario field; endogenous policing and crime
are outside the reference model.

\section{Workplaces, homes, and service venues}\label{app:buildings}
An actual vacancy additionally has a location, hourly wage quote, payroll
history, required skill vector, schedule, contract, and training access.
Effort, hazard, and autonomy in Table~\ref{tab:jobs} are adjustable
physical and organizational baselines. Outdoor share and noise are
features for tastes or an explicit health extension, not duplicate
unexplained effort penalties.

\begingroup
\footnotesize
\setlength{\tabcolsep}{2pt}
\renewcommand{\arraystretch}{1.15}
\begin{longtable}{@{}>{\raggedright\arraybackslash}p{0.43\textwidth}>{\raggedright\arraybackslash}p{0.09\textwidth}>{\raggedright\arraybackslash}p{0.09\textwidth}>{\raggedright\arraybackslash}p{0.1\textwidth}>{\raggedright\arraybackslash}p{0.12\textwidth}>{\raggedright\arraybackslash}p{0.08\textwidth}@{}}
\caption{Workplace baseline attributes.}\label{tab:jobs}\\
\toprule
\textbf{Workplace} & \textbf{Effort} & \textbf{Hazard} & \textbf{Autonomy} & \textbf{Outdoor share} & \textbf{h} \\
\midrule
\endfirsthead
\multicolumn{6}{l}{\textit{Table \thetable{} (continued)}}\\
\toprule
\textbf{Workplace} & \textbf{Effort} & \textbf{Hazard} & \textbf{Autonomy} & \textbf{Outdoor share} & \textbf{h} \\
\midrule
\endhead
\midrule
\multicolumn{6}{r}{\textit{Continued on next page}}\\
\endfoot
\bottomrule
\endlastfoot
Crop farm / garden & 0.55 & 0.15 & 0.45 & 0.90 & 8 \\
Fishing / forestry & 0.65 & 0.55 & 0.50 & 0.95 & 8 \\
Mine / quarry & 0.80 & 0.75 & 0.15 & 0.40 & 8 \\
Kiln / steel / cement & 0.75 & 0.60 & 0.15 & 0.20 & 8 \\
Mill / textile / paper & 0.50 & 0.30 & 0.20 & 0.10 & 8 \\
Food processing / bakery & 0.45 & 0.20 & 0.35 & 0.10 & 8 \\
Clothing / furniture / tools & 0.40 & 0.25 & 0.60 & 0.10 & 8 \\
Machine shop / utilities & 0.45 & 0.35 & 0.50 & 0.20 & 8 \\
Shop / hospitality & 0.45 & 0.10 & 0.35 & 0.10 & 8 \\
Freight / construction & 0.65 & 0.45 & 0.40 & 0.80 & 8 \\
Teacher / librarian & 0.35 & 0.05 & 0.65 & 0.05 & 6 \\
Clinical worker / pharmacist & 0.50 & 0.15 & 0.55 & 0.05 & 8 \\
Administration & 0.25 & 0.03 & 0.40 & 0.05 & 8 \\
\end{longtable}
\endgroup

A dwelling's rent is for the unit, not automatically per resident.
Privacy and thermal comfort are indexes in $[0,1]$. Area utility uses
adult-equivalents, whereas the hard occupancy limit counts persons.
Condition, weather, actual utility delivery, neighborhood exposure, and
travel time modify the realized services. A block uses either the eight
basic-unit layout or the five family-unit layout, not both.

\begingroup
\footnotesize
\setlength{\tabcolsep}{2pt}
\renewcommand{\arraystretch}{1.15}
\begin{longtable}{@{}>{\raggedright\arraybackslash}p{0.14\textwidth}>{\raggedright\arraybackslash}p{0.055\textwidth}>{\raggedright\arraybackslash}p{0.055\textwidth}>{\raggedright\arraybackslash}p{0.055\textwidth}>{\raggedright\arraybackslash}p{0.055\textwidth}>{\raggedright\arraybackslash}p{0.055\textwidth}>{\raggedright\arraybackslash}p{0.3\textwidth}>{\raggedright\arraybackslash}p{0.2\textwidth}@{}}
\caption{Residential types and construction requirements.}\label{tab:homes}\\
\toprule
\textbf{Dwelling} & \textbf{Area} & \textbf{Max. people} & \textbf{Shel- ter} & \textbf{Temp.} & \textbf{Priv.} & \textbf{Facilities} & \textbf{Package} \\
\midrule
\endfirsthead
\multicolumn{8}{l}{\textit{Table \thetable{} (continued)}}\\
\toprule
\textbf{Dwelling} & \textbf{Area} & \textbf{Max. people} & \textbf{Shel- ter} & \textbf{Temp.} & \textbf{Priv.} & \textbf{Facilities} & \textbf{Package} \\
\midrule
\endhead
\midrule
\multicolumn{8}{r}{\textit{Continued on next page}}\\
\endfoot
\bottomrule
\endlastfoot
Makeshift shelter & 12 m$^2$ & 2 & 0.45 & 0.30 & 0.45 & No private water, sanitation, or kitchen & 0.25 A \\
Worker dormitory & 80 m$^2$ & 8 & 0.85 & 0.50 & 0.20 & Shared water, washroom, kitchen & B \\
Basic apartment & 36 m$^2$ & 3 & 0.95 & 0.65 & 0.80 & Private kitchen/washroom; utility connections & One of 8 units in E \\
Family apartment & 60 m$^2$ & 4 & 1.00 & 0.80 & 0.90 & Private facilities; optional powered ventilation & One of 5 units in E \\
Cottage & 55 m$^2$ & 4 & 0.90 & 0.70 & 0.95 & Kitchen/washroom; garden; lower density & B \\
Comfortable house & 90 m$^2$ & 5 & 1.00 & 0.90 & 1.00 & Private facilities; garden; powered comfort option & C \\
\end{longtable}
\endgroup

Connected lighting uses 0.02 electricity units per occupied dwelling-day.
Optional powered ventilation uses 0.03 per adult-equivalent/day and can
add up to 0.10 thermal comfort when climatically useful, capped at one.
Without electricity the powered increment disappears, but structural
shelter does not. Garden, density, architectural style, quietness, and
known social ties are the small home-affinity features. Unhoused persons
have shelter and privacy zero and use only actually available public
facilities.

Venue quality is $q=c(0.6+0.4\bar s)u$, with condition $c$, relevant
mean skill $\bar s$, and supplied fraction $u$, all in $[0,1]$.
For complementary quality inputs, $u$ is their minimum supplied/required
ratio. Essential safety inputs are hard prerequisites. Capacity limits
produce queues or rejected bookings rather than fully served fictitious
visitors. Inputs listed below are per visit unless otherwise stated;
ordinary building maintenance and installed equipment are additional.

\begingroup
\footnotesize
\setlength{\tabcolsep}{2pt}
\renewcommand{\arraystretch}{1.15}
\begin{longtable}{@{}>{\raggedright\arraybackslash}p{0.14\textwidth}>{\raggedright\arraybackslash}p{0.23\textwidth}>{\raggedright\arraybackslash}p{0.25\textwidth}>{\raggedright\arraybackslash}p{0.1\textwidth}>{\raggedright\arraybackslash}p{0.2\textwidth}@{}}
\caption{Dated service technologies and household time.}\label{tab:venues}\\
\toprule
\textbf{Place} & \textbf{Building, staff, capacity} & \textbf{Variable inputs} & \textbf{Time} & \textbf{Delivered service} \\
\midrule
\endfirsthead
\multicolumn{5}{l}{\textit{Table \thetable{} (continued)}}\\
\toprule
\textbf{Place} & \textbf{Building, staff, capacity} & \textbf{Variable inputs} & \textbf{Time} & \textbf{Delivered service} \\
\midrule
\endhead
\midrule
\multicolumn{5}{r}{\textit{Continued on next page}}\\
\endfoot
\bottomrule
\endlastfoot
Market shop & C; 4 general workers; 80 transactions & 0.005 paper; inventory acquired separately & 0.10 h + queue & Goods and lower shopping/search time; no extra consumption utility \\
Communal kitchen & C + 1 machinery; 4 craft workers; 20 meal-equivalents & Food and cooking recipe inputs & 0.10 h + queue & Access to preparation for homes without kitchens \\
Restaurant & C + 1 machinery; 8 staff; 30 diners & 1 prepared meal + 0.02 water + 0.002 soap & 0.75 h + queue & Food nutrients/Q once; social +0.15q; recreation +0.10q \\
Bathhouse & C + 1 machinery; 4 staff; 40 visits & 0.5 water + 0.01 soap + 0.02 electricity & 0.40 h + queue & Hygiene for one day at access quality q \\
Park & Landscaped land + 0.25 A; 2 staff; 100 visitors & Maintenance labor; 0.01 water per visit-equivalent & 1 h & Recreation +0.5q; social +0.2q only when meeting others \\
Social club / tavern & C; 6 staff; 40 visits & 0.01 water + 0.002 soap; optional food bought separately & 1 h & Social +0.5q; recreation +0.3q \\
Theatre / cinema & E + 2 machinery; 8 staff; 100 seats/show & 0.002 electricity/seat + 0.002 paper & 1.5 h & Recreation +0.8q; reading/culture +0.2q \\
Library & C; 4 literate staff; 40 simultaneous users & Books as circulating stock; 0.005 electricity/user-hour & 1 h & Reading +0.5q times title novelty; books not consumed immediately \\
Religious venue & C; 3 staff; 80 attendees/service & 0.002 paper + 0.005 water & 1 h & Religion +1q; social +0.15q \\
Clinic & C + 2 machinery; 2 clinicians + 4 support; 12 visits & 0.2 medicines + 0.02 soap + 0.05 water + 0.02 electricity & 0.5 h + queue & Expected health improvement determined by condition and treatment \\
Hospital & E + 6 machinery; 4 clinicians + 8 support; 12 beds & Per occupied day: 0.5 medicines + 0.1 soap + 1 water + 0.2 electricity; food separately & Inpatient day & Recovery capacity; no direct enjoyment utility \\
Primary school & E; 4 teachers + 2 support; 80 pupils & Per pupil-day: 0.01 paper + 0.005 electricity; circulating books & 5 h/day & General education; supervision while in attendance \\
Vocational school & C + 3 machinery; 4 teachers; 32 students & Per student-day: 0.02 paper + 0.005 tools* + 0.02 electricity & 4 h/day & Occupational training \\
College & E + 4 machinery; 6 teachers + 2 support; 60 students & Per student-day: 0.03 paper + 0.03 electricity; books & 5 h/day & Advanced technical, clinical, or teaching study \\
Bus depot & D + 4 machinery; 8 drivers + 2 technicians; 4 vehicles & 0.02 electricity per vehicle-km; replacements from maintenance & Route-specific & Passenger-km; no direct consumption utility \\
Freight depot & D + 4 machinery; 8 drivers + 2 technicians; 4 vehicles & 0.03 electricity per vehicle-km; replacements from maintenance & No passenger time & Cargo-km; delivery and capacity constraints \\
\end{longtable}
\endgroup

The vehicle chassis are included in the depots' installed machinery;
vehicles are electric in this compressed technology set. Each bus has 20
passenger places and each freight vehicle carries 20 cargo-space units.
Use one cargo-space unit per ordinary commodity lot, two per machinery
module, and one per furniture kit; utilities use their networks instead.
Baseline walking and uncongested road speeds are 4 and 20 km/h.
Both outward and return journeys consume vehicle time. A route has a
scenario daily capacity, and traffic exceeding it queues; travel cannot
occur on an absent road. A zero-congestion scenario is the reference,
while congestion rules must be reported when enabled.

Unpaid alternatives remain feasible: an hour talking with an available
friend adds 0.35 social service; an hour walking adds 0.25 recreation;
an hour reading an owned book adds 0.5 reading service times novelty;
an hour of private religious practice adds 0.4 religious service.
Conversation requires both participants' time. A shop provides access
and convenience, not an additional utility reward for spending.

\section{Education and reference parameters}\label{app:parameters}
Skills are general, craft, technical, clinical, and teaching, each in
$[0,1]$. Trained positions require the relevant skill at least 0.5;
advanced positions require 0.8. The reference clinical, technical, and
teaching jobs additionally require the appropriate credential. General
jobs have no credential requirement. A course cannot lower an already
higher skill. Learning times below deliberately compress a game timescale;
they must not be mistaken for real educational durations.

\begingroup
\footnotesize
\setlength{\tabcolsep}{2pt}
\renewcommand{\arraystretch}{1.15}
\begin{longtable}{@{}>{\raggedright\arraybackslash}p{0.26\textwidth}>{\raggedright\arraybackslash}p{0.22\textwidth}>{\raggedright\arraybackslash}p{0.25\textwidth}>{\raggedright\arraybackslash}p{0.2\textwidth}@{}}
\caption{Course technologies.}\label{tab:education}\\
\toprule
\textbf{Program} & \textbf{Prerequisite} & \textbf{Full-attendance duration} & \textbf{Target skill} \\
\midrule
\endfirsthead
\multicolumn{4}{l}{\textit{Table \thetable{} (continued)}}\\
\toprule
\textbf{Program} & \textbf{Prerequisite} & \textbf{Full-attendance duration} & \textbf{Target skill} \\
\midrule
\endhead
\midrule
\multicolumn{4}{r}{\textit{Continued on next page}}\\
\endfoot
\bottomrule
\endlastfoot
Basic education & None & 360 school days & General 0.6 \\
Craft apprenticeship & General >=0.3 & 120 training days & Craft 0.6 \\
Advanced craft & Craft >=0.5 & 120 training days & Craft 0.85 \\
Technician training & General >=0.5 & 240 training days & Technical 0.7 \\
Clinical training & General >=0.6 & 360 training days & Clinical 0.7 \\
Teacher training & General >=0.6 & 300 training days & Teaching 0.7 \\
\end{longtable}
\endgroup

Teachers, classrooms, materials, and student time are reserved in the
same way as other service inputs. General education and professional
education use their catalogued daily schedules. Craft apprenticeship
uses four training hours/day at a vocational school or a funded employer
training place. Initial essential teachers and technicians are seeded
within the initial population; subsequent recruitment is internal.

Observed payroll reliability begins with a Beta$(9,1)$ prior and updates
one success or failure per due funded-day promise. School employment and
completion beliefs use Beta$(1,1)$ priors and observed individual
outcomes. Quoted daily output capacities are initial asset properties,
not monetary valuations. Registration information is visible only in
the relevant local market.

\begingroup
\footnotesize
\setlength{\tabcolsep}{2pt}
\renewcommand{\arraystretch}{1.15}
\begin{longtable}{@{}>{\raggedright\arraybackslash}p{0.31\textwidth}>{\raggedright\arraybackslash}p{0.24\textwidth}>{\raggedright\arraybackslash}p{0.38\textwidth}@{}}
\caption{Numerical defaults for reproducible experiments. Adjustment gains are hypotheses to test, not established stability bounds.}\label{tab:defaults}\\
\toprule
\textbf{Parameter} & \textbf{Value} & \textbf{Interpretation} \\
\midrule
\endfirsthead
\multicolumn{3}{l}{\textit{Table \thetable{} (continued)}}\\
\toprule
\textbf{Parameter} & \textbf{Value} & \textbf{Interpretation} \\
\midrule
\endhead
\midrule
\multicolumn{3}{r}{\textit{Continued on next page}}\\
\endfoot
\bottomrule
\endlastfoot
Household daily discount & 0.9999 & Applied to private forecast utility \\
Forecast exponential gain & 0.1 & Own observed quantities and income \\
Budget adjustment coefficient & 0.1 & Equation (budget rule) \\
Household reserve target & 7 days & Recent essential expenditure \\
Goods quote adjustment & 0.02 / daily round & Normalized pre-rationing imbalance \\
Wage offer adjustment & 0.02 / review & New offers only \\
Rent / tuition adjustment & 0.01 / review & New contracts only \\
Firm planning review & 7 days & Seeded staggered offset \\
Investment review & 30 days & At most one new module considered for execution \\
Employment / lease term & 31 / 30 days & Employment includes paid onboarding \\
Firm demand-share exploration & 0.05 & Capacity-weighted term in demand assignment \\
Inventory release gain & 0.25 & Deviation from physical buffer target \\
Inventory buffer & Up to 3 days & Limited by shelf life \\
Routine household horizon & One day; 30-day reference criterion & Fixed-choice LP; review uses rollouts \\
Move / job horizon & 90 days & Includes switching cost \\
Education evaluation horizon & Remaining course + 360 days & Local forecast only \\
Firm planning horizon & Maximum of 31, 7 + longest lag, and 1 + remaining term & Covers hiring commitment \\
Investment forecast horizon & 360 days & Includes construction and working capital \\
Terminal resale fraction & 0.5 & A belief applied to observed quotes, not a floor \\
Minimum reference food reserve & 35 days & An initial endowment, not replenishment \\
\end{longtable}
\endgroup

A complete scenario supplies the population and household partition,
map and network, asset ownership, initial cash, inventories and ages,
skills and credentials, natural stocks, latent fixed tastes, seed,
initial positive quotes, and public-policy schedule. For an initial
currency convention one may set every unit quote to one and distribute a
declared money stock; these initial quotes are not remembered targets.
The scenario also supplies posted jobs and leases or their bootstrap
matching phase. The article defines a parameterized transition algorithm,
not one privileged starting allocation or a demonstrated equilibrium.


\section{Notation guide}\label{app:notation}
Superscripts and subscripts are part of symbol definitions; bare letters
with several roles below should not be interchanged.
\begin{longtable}{@{}p{0.28\textwidth}p{0.66\textwidth}@{}}
\toprule Symbol & Meaning and domain\\\midrule
$\mathcal I,\HH,\mathcal F,\mathcal L,\mathcal G$ & Persons, households, establishments, locations and goods.\\
$\mathcal J_{a,t}$ & Agent information filtration (distinct from establishment set).\\
$i,h,f,g,k,t$ & Person, household, firm, good, skill/need and day indexes.\\
$P_i;\ P_i(o);\ P_e$ & Protein; bounded option taste bonus; course progress.\\
$\lambda_{h,t};\ \lambda\in\operatorname{spec}(J)$ & Marginal utility of cash; Jacobian eigenvalue.\\
$\rho_i;\ \rho(A)$ & Risk aversion; spectral radius.\\
$a_{kr};\ a_o;\ a_i;\ a_{kv};\ a$ & Recipe inputs (also vector $a_r$); option attributes; aptitude; service coefficient; generic agent/action index as specified locally.\\
$h_{i,t};\ h_i^{\rm paid};\ h\in\HH$ & Health; actual paid hours; household index.\\
$H(z);\ H_{ik};\ H_{jk,t}$ & Satiation function; effective worker hours; batch labor allocation.\\
$q;\ q_v;\ q_{fg};\ q^{\rm sold}_{fm}$ & Verifiable grade; venue quality; input order; realized sales quantity.\\
$r;\ r_j;\ r_t;\ r_a(y)$ & Recipe index; job hazard; repair fraction; benchmark resource use.\\
$r,K$ in the fish equation & Regeneration rate and carrying capacity only.\\
$b_r;\ b_g;\ b;\ \delta b$ & Recipe output vector; inventory buffer days; benchmark resource endowment or local cash budget argument; finite cash increment.\\
$m_{a,t};\ m=(\ell,g,q)$ & Cash balance; commodity-market identifier.\\
$I_g;\ c_f;\ v_j$ & Physical lot count; asset condition; batch progress.\\
$z_{ik};\ z_f$ & Normalized household need; ten-unit power operating blocks.\\
$\beta;\ T_f;\ \eta_m$ & Daily discount; firm forecast horizon; log-price adjustment gain.\\
\bottomrule
\end{longtable}

\section{Static Leontief feasibility diagnostic}\label{app:feasibility}
The supplementary script \texttt{feasibility\_check.py} solves a continuous
steady-flow relaxation of a 17-route subnetwork of the catalog. It includes
grain, legumes, vegetables, fish, water, timber, boards, charcoal, ore,
limestone, clay, steel, tools, machinery, bricks, cement and electricity.
Only raw foods are final consumption; prepared foods and the other goods
are disabled in this restricted experiment. Unit effective productivity,
fertility one, sufficient renewable resources, zero transport and spoilage,
and initially installed capital are assumed.

For batch rates $x_r\ge0$, food quantities $c_g\ge0$ and maximum daily
batch capacity $K_r$, install fractional facility counts $x_r/K_r$.
Here $K_r$ is the minimum of slot/lag capacity and worker-hours divided
by batch-hours; electricity has 12 operating blocks/day. Add each
facility's daily building repair bill $10^{-4}$ times its construction
package and operating machinery replacement $5\,10^{-4}$ times installed
modules to its input column, divided by $K_r$. Include repair labor and
construction-tool wear, $0.01/80$ tools per repair hour. One eighth of a
worker dormitory (package B) provides shelter index 0.85 per adult and
its maintenance is a fixed final requirement; this is basic shelter,
not the maximum shelter index of one.

With this maintenance-adjusted net-output matrix $B$, solve
\[
 \min_{x,c\ge0}\ \ell^\top x+
 0.4(c_{\rm grain}+0.8c_{\rm legumes}+c_{\rm fish})+0.002
\]
subject to $Bx\ge c+\text{cooking inputs}+\text{drinking water}
+\text{dormitory repair}$, $1\le N(c)\le1.1$, $P(c)\ge0.5$,
$M(c)\ge0.5$. Drinking water is 0.03 lots/adult-day.
Cooking consumes 0.03 charcoal and 0.01 water per cooked energy unit;
vegetables may be eaten raw. The constant 0.002 is dormitory repair
hours per adult-day. Electricity balances include all upstream production.

The solved objective is 3.593200 hours/adult-day, comprising 3.229011
effective production and repair hours and 0.364188 actual domestic
cooking hours. The food bundle is 0.525038 grain, zero legumes,
0.358118 vegetables and 0.385432 fish, with $(N,P,M)=(1,0.5,0.5)$.
Maximum linear-constraint violation is below $10^{-12}$ in the recorded
run. These are computed feasibility results, not simulation observations.
They show one resource-feasible restricted steady-flow configuration
under the stated assumptions. They neither prove monetary coordination
nor cover clothing, hygiene, education, other services, startup delays,
integer facilities, resource regeneration limits, commuting or inactive
dependents. Adding other food technologies could lower this restricted
minimum, while imposing those omitted requirements can raise labor needs.

\twocolumn
\balance
\begin{thebibliography}{99}
\bibitem{grimm2020} V. Grimm et al., The ODD protocol for describing
agent-based and other simulation models: A second update to improve
clarity, replication, and structural realism,
\emph{Journal of Artificial Societies and Social Simulation} \textbf{23}(2), 7 (2020),
\url{https://doi.org/10.18564/jasss.4259}.
\bibitem{hahn1962} F. H. Hahn and T. Negishi, A theorem on
non-t\^atonnement stability, \emph{Econometrica} \textbf{30}, 463--469 (1962),
\url{https://doi.org/10.2307/1909889}.
\bibitem{galeshapley1962} D. Gale and L. S. Shapley, College admissions
and the stability of marriage, \emph{American Mathematical Monthly}
\textbf{69}, 9--15 (1962), \url{https://doi.org/10.1080/00029890.1962.11989827}.
\bibitem{kelsocrawford1982} A. S. Kelso, Jr. and V. P. Crawford,
Job matching, coalition formation, and gross substitutes,
\emph{Econometrica} \textbf{50}, 1483--1504 (1982),
\url{https://econweb.ucsd.edu/~vcrawfor/KelsoCrawford82EMT.pdf}.
\bibitem{roth1984} A. E. Roth, The evolution of the labor market for
medical interns and residents: A case study in game theory,
\emph{Journal of Political Economy} \textbf{92}, 991--1016 (1984),
\url{https://doi.org/10.1086/261272}.
\bibitem{gintis2007} H. Gintis, The dynamics of general equilibrium,
\emph{Economic Journal} \textbf{117}, 1280--1309 (2007),
\url{https://doi.org/10.1111/j.1468-0297.2007.02083.x}.
\bibitem{dawid2014} H. Dawid, S. Gemkow, P. Harting, S. van der Hoog,
and M. Neugart, Agent-based macroeconomic modeling and policy analysis:
The Eurace@Unibi model, Bielefeld Working Papers in Economics and
Management No. 01-2014 (2014),
\url{https://doi.org/10.4119/unibi/2915476}.
\bibitem{scarf1960} H. E. Scarf, Some examples of global instability of
the competitive equilibrium, \emph{International Economic Review}
\textbf{1}, 157--172 (1960),
\url{https://doi.org/10.2307/2556215}.
\bibitem{werning2026} I. Werning and G. Lorenzoni, T\^atonnement and
price setting in general equilibrium, NBER Working Paper No. 35205
(2026), \url{https://doi.org/10.3386/w35205}.
\bibitem{boyd2004} S. Boyd and L. Vandenberghe,
\emph{Convex Optimization} (Cambridge University Press, Cambridge, 2004),
\url{https://doi.org/10.1017/CBO9780511804441}.
\bibitem{agrawal2022} A. Agrawal, S. Boyd, D. Narayanan,
F. Kazhamiaka, and M. Zaharia, Allocation of fungible resources via a
fast, scalable price discovery method, \emph{Mathematical Programming
Computation} \textbf{14}, 593--622 (2022),
\url{https://stanford.edu/~boyd/papers/resource_alloc.html}.
\end{thebibliography}
\end{document}
